%% 3-12 対数：けたを数える道具（第12回・組 V3-6）  % noruby
%% 流れ：電卓で「10を何乗すると2？」（やってみよう）→ ①けた＝10を何乗するか（常用対数）→ ②log_a b と底  % noruby
%%       → ③積→和・累乗→かけ算（3-11 の指数のきまりから）→ ④y=log_2 x は y=2^x を y=x で折り返す  % noruby
%%       → ⑤底の変え方 → ⑥片対数・両対数グラフ → ⑦デシベル。e と ln は使わない（5-14 と予告だけ）。  % noruby
\makeatletter
% 説明の段の見出し（丸数字＋太字の1行）  \bkthrmstep{番号}{見出し}  % noruby
\newcommand\bkthrmstep[2]{\par\addvspace{6pt}\Needspace{5\baselineskip}\noindent
  {\bkcircled[bkOrange]{#1}\enspace\sffamily\bfseries\color{bkOrange!75!black}#2}\par\nobreak\vspace{2pt}}
% 式を、まん中に1行で置く  \bkthrmline{…}  % noruby
\newcommand\bkthrmline[1]{\par\addvspace{3pt}{\centering #1\par}\addvspace{3pt}}
% 指数の形と対数の形を、矢印で行き来する図  \bkthrmswap{指数の形}{対数の形}  % noruby
%   \bkthrmswap{指数の形}{対数の形}{1 なら上に「指数の形」「対数の形」の札をつける}  % noruby
\newcommand\bkthrmswap[3]{\begin{tikzpicture}[x=1mm, y=1mm]
  \node[bk node, minimum width=36mm, minimum height=11mm, font=\normalsize] (p) at (0,0) {#1};
  \node[bk node, minimum width=36mm, minimum height=11mm, font=\normalsize] (q) at (62,0) {#2};
  \draw[bk link] ([yshift=1.8mm]p.east) -- ([yshift=1.8mm]q.west);
  \draw[bk link] ([yshift=-1.8mm]q.west) -- ([yshift=-1.8mm]p.east);
  \ifnum#3=1
  \node[font=\sffamily\footnotesize, anchor=south] at (0,7.5) {\R{指数}{し|すう}の\R{形}{かたち}};
  \node[font=\sffamily\footnotesize, anchor=south] at (62,7.5) {\R{対数}{たい|すう}の\R{形}{かたち}};
  \node[font=\sffamily\scriptsize, anchor=south] at (31,3) {なおす};
  \fi
\end{tikzpicture}}
% 「ここに入る数は？」の空らん（数式の中で使う）  \bkthrmbox  % noruby
\newcommand\bkthrmbox{\text{\fboxsep=0pt\fboxrule=0.5pt\fbox{\rule{0pt}{0.75em}\rule{0.6em}{0pt}}}}
% 表で目立たせるセルの色  % noruby
\colorlet{bkthrmsel}{bkOrange!22}
\makeatother

\bkchapter{3}{12}{\R{対数}{たい|すう}：けたを\R{数}{かぞ}える\R{道具}{どう|ぐ}}{\bklevelrange{2}{3}}{\R{対数}{たい|すう}：けたを\R{数}{かぞ}える\R{道具}{どう|ぐ}}
\bkmichishirube{2-7,2-8,3-11}{3-16,5-14,9-8}

\begin{goal}
\begin{itemize}[label={\color{\bkTT}●}, leftmargin=1.6\zw, itemsep=0pt, topsep=1pt]
\item \bklevel{2}\enspace 「10を\R{何乗}{なん|じょう}すると、その\R{数}{かず}になるか」でけたを\R{数}{かぞ}え、$\log_ab$ の\R{意味}{い|み}を\R{説明}{せつ|めい}できる。
\item \bklevel{3}\enspace \R{対数}{たい|すう}のきまり（かけ\R{算}{ざん}がたし\R{算}{ざん}になる）を、\R{指数}{し|すう}のきまりから\R{作}{つく}って\R{使}{つか}える。
\item \bklevel{3}\enspace \R{片対数}{かた|たい|すう}グラフやデシベルで、\R{倍々}{ばい|ばい}に\R{変}{か}わる\R{量}{りょう}を\R{直線}{ちょく|せん}やたし\R{算}{ざん}であつかえる。
\end{itemize}
\end{goal}

\begin{hatena}
\noindent\begin{minipage}[c]{0.6\linewidth}
\R{音}{おと}の\R{大}{おお}きさは、デシベル（dB）という\R{単位}{たん|い}で\R{表}{あらわ}します。ささやき\R{声}{ごえ}は\R{約}{やく}30\,dB、ふつうの\R{会話}{かい|わ}は\R{約}{やく}60\,dB、\R{電車}{でん|しゃ}が\R{通}{とお}るときのガード\R{下}{した}は\R{約}{やく}100\,dB です。

60\,dB の\R{会話}{かい|わ}の\R{音}{おと}の\R{強}{つよ}さは、30\,dB のささやき\R{声}{ごえ}の2\R{倍}{ばい}かな？　じつは、もっとずっと\R{大}{おお}きくちがいます。\R{何倍}{なん|ばい}ちがうのでしょう。
\end{minipage}\hfill
\begin{minipage}[c]{0.37\linewidth}\centering
\begin{tikzpicture}[x=1mm, y=1mm]
  \fill[white, rounded corners=1mm] (0,0) rectangle (9,48);
  \shade[bottom color=bkGreen!40, top color=bkRed!55, rounded corners=1mm] (0.8,0.8) rectangle (8.2,47.2);
  \draw[bkInk!70, line width=0.7pt, rounded corners=1mm] (0,0) rectangle (9,48);
  \foreach \d in {0,20,...,120} {\draw[bkInk, line width=0.5pt] (9,\d*0.4) -- (10.8,\d*0.4);}
  \foreach \d in {0,40,80,120} {\node[font=\sffamily\scriptsize, anchor=east] at (-1.5,\d*0.4) {\d};}
  \node[font=\sffamily\scriptsize, anchor=south] at (4.5,49.5) {dB};
  \foreach \d/\tx in {30/ささやき\R{声}{ごえ}, 60/ふつうの\R{会話}{かい|わ}, 100/ガード\R{下}{した}} {
    \fill[bkInk] (9,\d*0.4) circle (0.9);
    \draw[bkInk!70, line width=0.5pt] (9,\d*0.4) -- (13,\d*0.4);
    \node[font=\sffamily\scriptsize, anchor=west] at (13.5,\d*0.4) {\d\ \tx};}
\end{tikzpicture}
\end{minipage}
\end{hatena}

この\R{章}{しょう}では、「10を\R{何回}{なん|かい}かけたか」で\R{数}{かず}の\R{大}{おお}きさを\R{表}{あらわ}す\R{道具}{どう|ぐ}、\textbf{\R{対数}{たい|すう}}を\R{作}{つく}ります。3-11の\R{指数}{し|すう}を、\R{反対}{はん|たい}から\R{見}{み}る\R{道具}{どう|ぐ}です。

%% ---------------- やってみよう・思い出そう ----------------
\begin{yattemiyou}[10を\R{何乗}{なん|じょう}すると2になる？]
\R{電卓}{でん|たく}の \texttt{x\^{}y} のキー（\R{表計算}{ひょう|けい|さん}なら \texttt{=10\^{}0.5} のような\R{式}{しき}）を\R{使}{つか}います。
\begin{enumerate}[label={\bkcircled[bkGreen]{\arabic*}}, leftmargin=2.4\zw, itemsep=1pt, topsep=3pt]
\item $10^0=1$、$10^1=10$ です。2は1と10のあいだなので、\R{指数}{し|すう}は0と1のあいだにありそうです。
\item $10^{0.5}$ を\R{計算}{けい|さん}しよう。2より\R{大}{おお}きいかな、\R{小}{ちい}さいかな？
\item \R{指数}{し|すう}を0.4、0.3、0.31……と\R{変}{か}えて、2にいちばん\R{近}{ちか}くなる\R{指数}{し|すう}を、\R{小数第}{しょう|すう|だい}3\R{位}{い}までさがそう。
\item \R{同}{おな}じように、10を\R{何乗}{なん|じょう}すると3になるかもさがそう。
\end{enumerate}
\end{yattemiyou}

\noindent \R{結果}{けっ|か}は\R{次}{つぎ}のようになります。2-8で、2\R{乗}{じょう}ではさんで $\sqrt{2}$ に\R{近}{ちか}づいたのと\R{同}{おな}じやり\R{方}{かた}です。

\begin{bkcenter}
\begin{bktable}{colspec={Q[l,m,wd=20mm] *{6}{Q[c,m,wd=14mm]}}, rowsep=1.5pt, cells={font=\small}, row{1}={bg=bkTheme!10}, cell{2}{6}={bg=bkthrmsel}}
\R{指数}{し|すう} $x$ & 0.5 & 0.4 & 0.3 & 0.31 & 0.301 & 0.302 \\
$10^x$ & 3.162 & 2.512 & 1.995 & 2.042 & 1.9999 & 2.004 \\
\end{bktable}
\bktabcaption{$10^{0.3}$ は2より\R{少}{すこ}し\R{小}{ちい}さく、$10^{0.31}$ は\R{少}{すこ}し\R{大}{おお}きい。$10^{0.301}$ がいちばん2に\R{近}{ちか}い。}
\end{bkcenter}

\noindent 2は、10の\R{約}{やく}0.301\R{乗}{じょう}です。\R{同}{おな}じようにさがすと、3は10の\R{約}{やく}0.477\R{乗}{じょう}です（$10^{0.477}=2.999\cdots$）。

\begin{omoidasou}
\bkrecall{2-7}{$1000=10^3$（0が3つ）、\ $0.01=10^{-2}$。10\R{倍}{ばい}するごとに、10の\R{指数}{し|すう}が1ふえる。}
\bkrecall{2-8}{2\R{乗}{じょう}して2に\R{近}{ちか}い\R{数}{かず}を、はさみながら\R{少}{すこ}しずつさがした（$\sqrt{2}=1.414\cdots$）。}
\bkrecall{3-11}{$a^{x+h}=a^x\times a^h$、\ $(a^m)^n=a^{mn}$、\ $a^0=1$、\ $a^{-x}=\dfrac{1}{a^x}$}
\bkrecall{3-11}{$x^{\frac{1}{2}}=\sqrt{x}$。\R{指数関数}{し|すう|かん|すう} $y=2^x$ のグラフは\R{右上}{みぎ|あ}がりで、$y$ はいつも\R{正}{せい}。}
\end{omoidasou}


\clearpage
%% ---------------- 数学の道具 ----------------
\bksection{doguu}{\R{数学}{すう|がく}の\R{道具}{どう|ぐ}}

\bksubsection{1}{けたを\R{数}{かぞ}える：10を\R{何乗}{なん|じょう}するか}

\leavevmode\label{bk:3-12:keta}\R{大}{おお}きな\R{数}{かず}や\R{小}{ちい}さな\R{数}{かず}は、10を\R{何乗}{なん|じょう}したものかで\R{表}{あらわ}すと、\R{大}{おお}きさがすぐにわかります（2-7）。

\begin{bkcenter}
\begin{bktable}{colspec={Q[l,m,wd=34mm] *{7}{Q[c,m,wd=11mm]}}, rowsep=1.5pt, colsep=3pt, cells={font=\small}, column{1}={font=\sffamily\small}, row{3}={bg=bkthrmsel}}
\R{数}{かず} & 0.01 & 0.1 & 1 & 10 & 100 & 1000 & 10000 \\
10の\R{累乗}{るい|じょう}で & $10^{-2}$ & $10^{-1}$ & $10^0$ & $10^1$ & $10^2$ & $10^3$ & $10^4$ \\
10を\R{何乗}{なん|じょう}するか & $-2$ & $-1$ & 0 & 1 & 2 & 3 & 4 \\
\end{bktable}
\bktabcaption{10\R{倍}{ばい}するごとに、「10を\R{何乗}{なん|じょう}するか」は1ずつふえる。}
\end{bkcenter}

1000は0が3つで4けた、10000は0が4つで5けたです。「10を\R{何乗}{なん|じょう}するか」は、\R{整数}{せい|すう}なら\R{数}{かず}のけた\R{数}{すう}より1\R{小}{ちい}さい\R{数}{かず}で、けたを\R{数}{かぞ}えていることになります。

では、2や3や2560のように、10の\R{累乗}{るい|じょう}でない\R{数}{かず}はどうでしょう。やってみようで、2は10の\R{約}{やく}0.301\R{乗}{じょう}、3は10の\R{約}{やく}0.477\R{乗}{じょう}とわかりました。\R{図}{ず}\ref{bk:3-12:fig-keta}のように、2\R{本}{ほん}の\R{目}{め}もりを\R{並}{なら}べるとわかりやすくなります。

\begin{bkfigure}
\begin{tikzpicture}[x=1mm, y=1mm]
  \draw[bkInk, line width=0.8pt] (0,0) -- (108,0);
  \foreach \k in {0,...,3} {\draw[bkInk, line width=0.8pt] (\k*34,-1.8) -- (\k*34,1.8);
    \node[font=\small, anchor=north] at (\k*34,-2.6) {\k};}
  \foreach \k/\v in {0/1,1/10,2/100,3/1000} {\node[font=\small, anchor=south] at (\k*34,2.6) {\v};}
  % 1と10のあいだ（2〜9）・10と100のあいだ（20〜90）の目もり  % noruby
  \foreach \m in {2,...,9} {\pgfmathsetmacro\xx{34*log10(\m)}\draw[bkBlue, line width=0.6pt] (\xx,0) -- (\xx,1.3);
    \pgfmathsetmacro\xy{34+34*log10(\m)}\draw[bkGray, line width=0.5pt] (\xy,0) -- (\xy,1.1);
    \pgfmathsetmacro\xz{68+34*log10(\m)}\draw[bkGray, line width=0.5pt] (\xz,0) -- (\xz,1.1);}
  \pgfmathsetmacro\xa{34*log10(2)}\pgfmathsetmacro\xb{34*log10(3)}
  \draw[bkBlue!80!black, line width=0.5pt] (\xa,1.3) -- (\xa,10.5);
  \draw[bkBlue!80!black, line width=0.5pt] (\xb,1.3) -- (\xb,10.5);
  \node[font=\footnotesize, text=bkBlue!80!black, anchor=south east, inner sep=1pt] at (\xa,10.8) {2};
  \node[font=\footnotesize, text=bkBlue!80!black, anchor=south west, inner sep=1pt] at (\xb,10.8) {3};
  \draw[bkBlue!80!black, line width=0.5pt, dashed] (\xa,0) -- (\xa,-8.5);
  \draw[bkBlue!80!black, line width=0.5pt, dashed] (\xb,0) -- (\xb,-8.5);
  \node[font=\footnotesize, text=bkBlue!80!black, anchor=north east, inner sep=1pt] at (\xa,-8.8) {0.301};
  \node[font=\footnotesize, text=bkBlue!80!black, anchor=north west, inner sep=1pt] at (\xb,-8.8) {0.477};
  \node[font=\sffamily\footnotesize, anchor=east] at (-3,6) {\R{数}{かず}};
  \node[font=\sffamily\footnotesize, anchor=east] at (-3,-6) {10の\R{何乗}{なん|じょう}};
  \foreach \k in {0,...,2} {\draw[-{Stealth[length=3pt]}, bkOrange, line width=0.7pt] (\k*34+5,-16) -- (\k*34+29,-16);
    \node[font=\sffamily\scriptsize, text=bkOrange!80!black, anchor=north] at (\k*34+17,-17) {×10};}
\end{tikzpicture}
\bkfigcaption{\R{上}{うえ}は\R{数}{かず}、\R{下}{した}は「10の\R{何乗}{なん|じょう}か」。10\R{倍}{ばい}するごとに、\R{同}{おな}じ\R{長}{なが}さ（1\R{目}{め}もり）だけ\R{右}{みぎ}へ\R{進}{すす}む。}\label{bk:3-12:fig-keta}
\end{bkfigure}

\R{上}{うえ}の\R{目}{め}もりは、10\R{倍}{ばい}するごとに\R{同}{おな}じ\R{長}{なが}さだけ\R{進}{すす}みます。かけ\R{算}{ざん}が、\R{長}{なが}さのたし\R{算}{ざん}になっているのです。この\R{目}{め}もりを\textbf{\R{対数目}{たい|すう|め}もり}といいます。

\bkthrmstep{1}{\R{名前}{な|まえ}と\R{記号}{き|ごう}をつける}
「10を\R{何乗}{なん|じょう}すると $b$ になるか」を\R{表}{あらわ}す\R{数}{かず}を、$\log_{10}b$ と\R{書}{か}き、「ログ・\R{十}{じゅう}・ビー」と\R{読}{よ}みます。たとえば、
\[ \log_{10}1000=3,\qquad \log_{10}0.01=-2 \]
\[ \log_{10}2=0.301\cdots,\qquad \log_{10}3=0.477\cdots \]
このような\R{数}{かず}を\textbf{\R{常用対数}{じょう|よう|たい|すう}}といいます。\bkidx{じょうようたいすう}{\R{常用対数}{じょう|よう|たい|すう}}\R{常用対数}{じょう|よう|たい|すう}は、10を\R{省}{はぶ}いて $\log2$ と\R{書}{か}くことも\R{多}{おお}く、\R{電卓}{でん|たく}の \texttt{log} のキー（\R{表計算}{ひょう|けい|さん}なら \texttt{=LOG10(2)}）で\R{求}{もと}められます。

\bkthrmstep{2}{けたの\R{数}{かず}がわかる}
2560のけたを\R{考}{かんが}えます。1000＜2560＜10000なので、\R{図}{ず}\ref{bk:3-12:fig-keta}の\R{目}{め}もりでは3と4のあいだです。\R{電卓}{でん|たく}で $\log_{10}2560=3.408\cdots$ です。
\begin{enumerate}[label={\bkcircled[bkTheme]{\arabic*}}, leftmargin=2.4\zw, itemsep=1pt, topsep=2pt]
\item $\log_{10}b$ が3\R{以上}{い|じょう}4\R{未満}{み|まん}なら、$b$ は $10^3=1000$ \R{以上}{い|じょう}、$10^4=10000$ \R{未満}{み|まん}。
\item 1000\R{以上}{い|じょう}10000\R{未満}{み|まん}の\R{整数}{せい|すう}は、4けた。
\item つまり、（$\log_{10}b$ の\R{整数}{せい|すう}の\R{部分}{ぶ|ぶん}）＋1 が、$b$ のけたの\R{数}{かず}です。
\end{enumerate}
\R{常用対数}{じょう|よう|たい|すう}は、まさに「けたを\R{数}{かぞ}える\R{道具}{どう|ぐ}」です。

\bksubsection{2}{$\log_ab$：$a$ を\R{何乗}{なん|じょう}すると $b$ になるか}

\leavevmode\label{bk:3-12:log}10のかわりに、ほかの\R{数}{かず}を\R{使}{つか}うこともできます。3-11の\R{細菌}{さい|きん}は、20\R{分}{ぷん}ごとに2\R{倍}{ばい}になりました。32\R{倍}{ばい}になるのは\R{何回}{なん|かい}\R{分}{わ}かれたときでしょう。$2^5=32$ なので、5\R{回}{かい}です。「2を\R{何乗}{なん|じょう}すると32になるか」の\R{答}{こた}えの5を、
\[ \log_2 32=5 \]
と\R{書}{か}きます。\R{一般}{いっ|ぱん}に、「$a$ を\R{何乗}{なん|じょう}すると $b$ になるか」を\R{表}{あらわ}す\R{数}{かず}を $\log_ab$ と\R{書}{か}き、$a$ を\textbf{\R{底}{てい}}といいます。\bkidx{てい}{\R{底}{てい}}このような\R{数}{かず}を、まとめて\textbf{\R{対数}{たい|すう}}といいます。\bkidx{たいすう}{\R{対数}{たい|すう}}3-11の\R{指数関数}{し|すう|かん|すう} $y=a^x$ の $a$ も、\R{底}{てい}とよびます。

\begin{bkfigure}
\bkthrmswap{$2^5=32$}{$\log_2 32=5$}{1}\par\vspace{4pt}
\bkthrmswap{$a^p=b$}{$\log_ab=p$}{0}
\bkfigcaption{\R{同}{おな}じことを、2つの\R{形}{かたち}で\R{書}{か}く。\R{指数}{し|すう}の\R{形}{かたち}の\R{右肩}{みぎ|かた}の\R{数}{かず}（\R{指数}{し|すう}）が、\R{対数}{たい|すう}の\R{答}{こた}えになる。}\label{bk:3-12:fig-swap}
\end{bkfigure}

\begin{bkcenter}
\begin{bktable}{colspec={Q[l,m,wd=26mm] Q[l,m,wd=58mm] Q[c,m,wd=22mm]}, rowsep=1.5pt, cells={font=\small}, row{1}={font=\sffamily\bfseries\small, bg=bkTheme!10}}
\R{対数}{たい|すう} & \R{考}{かんが}え\R{方}{かた}（\R{指数}{し|すう}の\R{形}{かたち}） & \R{答}{こた}え \\
$\log_2 8$ & $2^{\bkthrmbox}=8$。$2\times2\times2=8$ & 3 \\
$\log_3 9$ & $3^{\bkthrmbox}=9$。$3\times3=9$ & 2 \\
$\log_{10}0.01$ & $10^{\bkthrmbox}=0.01=\dfrac{1}{100}$ & $-2$ \\
$\log_2\dfrac{1}{2}$ & $2^{\bkthrmbox}=\dfrac{1}{2}$。$2^{-1}=\dfrac{1}{2}$（3-11） & $-1$ \\
$\log_4 2$ & $4^{\bkthrmbox}=2$。$4^{\frac{1}{2}}=\sqrt{4}=2$（3-11） & $\dfrac{1}{2}$ \\
$\log_5 1$ & $5^{\bkthrmbox}=1$。$5^0=1$ & 0 \\
$\log_7 7$ & $7^{\bkthrmbox}=7$。$7^1=7$ & 1 \\
\end{bktable}
\bktabcaption{\R{対数}{たい|すう}は、\R{指数}{し|すう}の\R{形}{かたち}になおして、$\bkthrmbox$ に\R{入}{はい}る\R{数}{かず}をさがせば\R{求}{もと}められる。}
\end{bkcenter}

\R{最後}{さい|ご}の2つから、どんな\R{底}{てい} $a$ でも $\log_a1=0$、$\log_aa=1$ です。0\R{乗}{じょう}すると1、1\R{乗}{じょう}するとそのままだからです。

\bkthrmstep{1}{$a^{\log_ab}=b$}
$\log_ab$ は「$a$ を\R{何乗}{なん|じょう}すると $b$ になるか」の\R{答}{こた}えでした。だから、$a$ を $\log_ab$ \R{乗}{じょう}すると、ちゃんと $b$ になります。
\bkthrmline{\fbox{$a^{\log_ab}=b$}\qquad たとえば\quad $2^{\log_2 32}=2^5=32$,\quad $10^{\log_{10}2}=10^{0.301}=2$}
\R{当}{あ}たり\R{前}{まえ}に\R{見}{み}えますが、⑥で、グラフの\R{傾}{かたむ}きから\R{倍率}{ばい|りつ}を\R{読}{よ}むときに、この\R{見方}{み|かた}を\R{使}{つか}います。

\bkthrmstep{2}{\R{対数}{たい|すう}にできる\R{数}{かず}・できない\R{数}{かず}}
3-11で\R{見}{み}たように、\R{底}{てい} $a$ が\R{正}{せい}の\R{数}{かず}なら、$a^p$ は $p$ がどんな\R{数}{かず}でも\R{正}{せい}の\R{数}{かず}です。だから、$\log_a0$ や $\log_a(-4)$ のように、0やマイナスの\R{数}{かず}の\R{対数}{たい|すう}はありません。2を\R{何乗}{なん|じょう}しても、0やマイナスにはならないからです。また、\R{底}{てい}は1でない\R{正}{せい}の\R{数}{かず}にします。1は\R{何乗}{なん|じょう}しても1なので、「1を\R{何乗}{なん|じょう}すると5か」には\R{答}{こた}えがないからです。

\noindent\begin{minipage}[t]{0.42\linewidth}\vspace{0pt}
$\log$ は、\R{対数}{たい|すう}を\R{表}{あらわ}す\R{英語}{えい|ご} logarithm（ロガリズム）の\R{頭}{あたま}の3\R{文字}{も|じ}です。\R{底}{てい} $a$ は、log の\R{右下}{みぎ|した}に\R{小}{ちい}さく\R{書}{か}きます。
\end{minipage}\hfill
\begin{minipage}[t]{67mm}\vspace{0pt}\bkkigocard{$\log_ab$}{ログ・エー・ビー}{$a$ を\R{何乗}{なん|じょう}すると $b$ になるかを\R{表}{あらわ}す\R{数}{かず}（\R{対数}{たい|すう}）。$a$ が\R{底}{てい}}{logarithm}{3-12}\end{minipage}
\par\vspace{6pt}

\bksubsection{3}{かけ\R{算}{ざん}が、たし\R{算}{ざん}になる}

\leavevmode\label{bk:3-12:kimari}2の\R{累乗}{るい|じょう}と、その\R{対数}{たい|すう} $\log_2$ を\R{並}{なら}べてみます。
\begin{bkcenter}
\begin{bktable}{colspec={Q[l,m,wd=17mm] *{11}{Q[c,m,wd=8.6mm]}}, rowsep=1.5pt, colsep=1.5pt, cells={font=\small}, row{1}={bg=bkTheme!10}, cell{1}{4,5}={bg=bkthrmsel}, cell{2}{4,5}={bg=bkthrmsel}, cell{1}{7}={bg=bkBlue!15}, cell{2}{7}={bg=bkBlue!15}}
$b$ & 1 & 2 & 4 & 8 & 16 & 32 & 64 & 128 & 256 & 512 & 1024 \\
$\log_2b$ & 0 & 1 & 2 & 3 & 4 & 5 & 6 & 7 & 8 & 9 & 10 \\
\end{bktable}
\bktabcaption{$4\times8=32$ のとき、\R{対数}{たい|すう}は $2+3=5$。}
\end{bkcenter}
$4\times8=32$ です。それぞれの\R{対数}{たい|すう}を\R{見}{み}ると、4は2、8は3、32は5です。$2+3=5$ になっています。$8\times16=128$ でも、$3+4=7$ です。\textbf{かけ\R{算}{ざん}の\R{答}{こた}えの\R{対数}{たい|すう}は、\R{対数}{たい|すう}どうしのたし\R{算}{ざん}}になっています。

\bkthrmstep{1}{なぜたし\R{算}{ざん}になるのか（\R{指数}{し|すう}のきまりから\R{作}{つく}る）}
$\log_ab=p$、$\log_ac=q$ とします。\R{指数}{し|すう}の\R{形}{かたち}になおすと、$b=a^p$、$c=a^q$ です（\R{図}{ず}\ref{bk:3-12:fig-swap}）。
\begin{enumerate}[label={\bkcircled[bkTheme]{\arabic*}}, leftmargin=2.4\zw, itemsep=1pt, topsep=2pt]
\item かけると、$bc=a^p\times a^q$。
\item 3-11の $a^{x+h}=a^x\times a^h$ で、$a^p\times a^q=a^{p+q}$。だから $bc=a^{p+q}$。
\item これを\R{対数}{たい|すう}の\R{形}{かたち}になおすと、$\log_a(bc)=p+q$。$p$ と $q$ をもとにもどすと、
\end{enumerate}
\bkthrmline{\fbox{$\log_a(bc)=\log_ab+\log_ac$}}
「かけ\R{算}{ざん}は\R{指数}{し|すう}をたす」（3-11）を、\R{対数}{たい|すう}のことばで\R{言}{い}いなおしたものです。

\textbf{\R{電卓}{でん|たく}で\R{確}{たし}かめる}\quad \R{常用対数}{じょう|よう|たい|すう}でも\R{成}{な}り\R{立}{た}ちます。$6=2\times3$ なので、
\[ \log_{10}6=\log_{10}2+\log_{10}3=0.301+0.477=0.778 \]
\R{電卓}{でん|たく}で $\log_{10}6$ を\R{計算}{けい|さん}すると0.778151… です。たしかに\R{合}{あ}っています。

\bkthrmstep{2}{\R{累乗}{るい|じょう}は、かけ\R{算}{ざん}になる}
$b^3=b\times b\times b$ なので、①のきまりを2\R{回}{かい}\R{使}{つか}うと、
\[ \log_a(b^3)=\log_ab+\log_ab+\log_ab=3\log_ab \]
です。3でなくても\R{同}{おな}じです。$\log_ab=p$ とすると $b=a^p$ で、3-11の $(a^m)^n=a^{mn}$ から、
\[ b^r=(a^p)^r=a^{pr} \]
なので、$\log_a(b^r)=pr=r\log_ab$ です。$r$ は\R{分数}{ぶん|すう}やマイナスでもかまいません。
\bkthrmline{\fbox{$\log_ab^r=r\log_ab$}}
たとえば、$\log_{10}2^{10}=10\times\log_{10}2=10\times0.301=3.01$ です。$2^{10}$ は10の\R{約}{やく}3.01\R{乗}{じょう}で、1000より\R{少}{すこ}し\R{大}{おお}きい\R{数}{かず}です。たしかに $2^{10}=1024$ でした（2-7）。

\bkthrmstep{3}{$r=-1$ のとき：\R{逆数}{ぎゃく|すう}の\R{対数}{たい|すう}と、わり\R{算}{ざん}}
$\dfrac{1}{2}=2^{-1}$（3-11）なので、$r=-1$ として
\[ \log_{10}\frac{1}{2}=\log_{10}2^{-1}=-1\times\log_{10}2=-0.301 \]
です。\R{確}{たし}かめると、$10^{-0.301}=\dfrac{1}{10^{0.301}}=\dfrac{1}{2}$ です。わり\R{算}{ざん}は「\R{逆数}{ぎゃく|すう}をかける」ことなので（2-4）、①と\R{合}{あ}わせて、
\[ \log_a\frac{b}{c}=\log_a\Bigl(b\times\frac{1}{c}\Bigr)=\log_ab+\log_a\frac{1}{c}=\log_ab-\log_ac \]
たとえば、$\log_{10}5=\log_{10}\dfrac{10}{2}=\log_{10}10-\log_{10}2=1-0.301=0.699$ です。

\begin{column}{\R{計算尺}{けい|さん|じゃく}：\R{長}{なが}さをたして、かけ\R{算}{ざん}をする}
\R{電卓}{でん|たく}がなかったころ、\R{技術者}{ぎ|じゅつ|しゃ}は\R{計算尺}{けい|さん|じゃく}というものさしでかけ\R{算}{ざん}をしました。\R{図}{ず}\ref{bk:3-12:fig-keta}の\R{対数目}{たい|すう|め}もりを2\R{本}{ほん}\R{用意}{よう|い}し、ずらして\R{長}{なが}さをたします。$2\times3$ なら、
\begin{enumerate}[label={\bkcircled[bkNavy]{\arabic*}}, leftmargin=2.4\zw, itemsep=0pt, topsep=2pt]
\item \R{下}{した}のものさしの1を、\R{上}{うえ}のものさしの2に\R{合}{あ}わせる（$\log2$ の\R{長}{なが}さだけずらす）。
\item \R{下}{した}のものさしの3の\R{真上}{ま|うえ}を\R{読}{よ}む。$\log2+\log3=\log6$ なので、6の\R{目}{め}もりがある。
\end{enumerate}
\par\vspace{2pt}
\centering
\begin{tikzpicture}[x=1mm, y=1mm]
  \filldraw[fill=bkYellowBg, draw=bkYellow!80!black, line width=0.6pt] (-3,6) rectangle (93,14);
  \filldraw[fill=white, draw=bkGray, line width=0.6pt] (24.1,-2) rectangle (93,6);
  \foreach \m in {1,...,9} {\pgfmathsetmacro\xx{80*log10(\m)}\draw[bkInk, line width=0.5pt] (\xx,6) -- (\xx,8.5);
    \node[font=\footnotesize, anchor=south, inner sep=0.8pt] at (\xx,8.6) {\m};}
  \draw[bkInk, line width=0.5pt] (80,6) -- (80,8.5); \node[font=\footnotesize, anchor=south, inner sep=0.8pt] at (80,8.6) {10};
  \foreach \m in {1,...,3} {\pgfmathsetmacro\xx{24.08+80*log10(\m)}\draw[bkInk, line width=0.5pt] (\xx,6) -- (\xx,3.5);
    \node[font=\footnotesize, anchor=north, inner sep=0.8pt] at (\xx,3.4) {\m};}
  \draw[bkRed, line width=0.9pt] (62.25,-3) -- (62.25,8.3);
  \draw[{Stealth[length=3pt]}-{Stealth[length=3pt]}, bkBlue, line width=0.7pt] (0,-6) -- (24.08,-6);
  \node[font=\sffamily\scriptsize, text=bkBlue!80!black, anchor=north] at (12,-6.8) {log 2 の\R{長}{なが}さ};
  \draw[{Stealth[length=3pt]}-{Stealth[length=3pt]}, bkOrange, line width=0.7pt] (24.08,-6) -- (62.25,-6);
  \node[font=\sffamily\scriptsize, text=bkOrange!80!black, anchor=north] at (43.2,-6.8) {log 3 の\R{長}{なが}さ};
  \node[font=\sffamily\scriptsize, text=bkRed!80!black, anchor=west] at (63.5,-4.5) {6を\R{読}{よ}む};
\end{tikzpicture}
\end{column}

\bksubsection{4}{$y=\log_2x$ のグラフ：\R{指数関数}{し|すう|かん|すう}を\R{折}{お}り\R{返}{かえ}す}

\leavevmode\label{bk:3-12:graph}$y=\log_2x$ は、$x$ を\R{決}{き}めると $y$ が1つ\R{決}{き}まる\R{関数}{かん|すう}です（3-3）。\R{表}{ひょう}を\R{作}{つく}ります。\R{対数}{たい|すう}の\R{形}{かたち}を\R{指数}{し|すう}の\R{形}{かたち}になおすと、$y=\log_2x$ は $x=2^y$ と\R{同}{おな}じです。

\begin{bkcenter}
\begin{bktable}{colspec={Q[l,m,wd=24mm] *{6}{Q[c,m,wd=12mm]}}, rowsep=1.5pt, cells={font=\small}}
$x$ & $\dfrac{1}{4}$ & $\dfrac{1}{2}$ & 1 & 2 & 4 & 8 \\
$y=\log_2x$ & $-2$ & $-1$ & 0 & 1 & 2 & 3 \\
\end{bktable}
\end{bkcenter}

3-11の $y=2^x$ の\R{表}{ひょう}（$x=-2$、$-1$、0、1、2、3 で $y=\dfrac{1}{4}$、$\dfrac{1}{2}$、1、2、4、8）とくらべると、$x$ と $y$ の\R{行}{ぎょう}が\R{入}{い}れかわっています。だから、$y=2^x$ のグラフの\R{点}{てん} $(2,\ 4)$ は、$y=\log_2x$ では\R{点}{てん} $(4,\ 2)$ になります。

\noindent\begin{minipage}[c]{0.56\linewidth}\centering
\begin{tikzpicture}
\begin{axis}[bkplot, width=78mm, height=78mm, xmin=-2.6, xmax=8.8, ymin=-2.6, ymax=8.8, xtick={-2,2,4,6,8}, ytick={-2,2,4,6,8}, xlabel={$x$}, ylabel={$y$}, tick label style={fill=white, inner sep=1pt}]
  \addplot[bkthin, dashed, domain=-2.5:8.3] {x};
  \addplot[bkcurve2, domain=-2.5:3.1] {2^x};
  \addplot[bkcurve, domain=0.17:8.5] {ln(x)/ln(2)};
  \foreach \p/\q in {1/2,2/4,3/8} {\edef\tmp{\noexpand\draw[bkGray, line width=0.5pt, densely dotted] (axis cs:\p,\q) -- (axis cs:\q,\p);}\tmp}
  \addplot[only marks, mark=*, mark size=1.6pt, bkOrange!80!black] coordinates {(0,1) (1,2) (2,4) (3,8)};
  \addplot[only marks, mark=*, mark size=1.6pt, bkBlue!80!black] coordinates {(1,0) (2,1) (4,2) (8,3)};
  \node[font=\sffamily\footnotesize, text=bkOrange!80!black, anchor=east, fill=white, inner sep=1pt] at (axis cs:2.75,7.4) {$y=2^x$};
  \node[font=\sffamily\footnotesize, text=bkBlue!80!black, anchor=north, fill=white, inner sep=1pt] at (axis cs:6.4,1.5) {$y=\log_2x$};
  \node[font=\sffamily\footnotesize, text=bkGray!60!black, anchor=north west, fill=white, inner sep=1pt] at (axis cs:7.9,7.4) {$y=x$};
\end{axis}
\end{tikzpicture}
\end{minipage}\hfill
\begin{minipage}[c]{0.41\linewidth}
オレンジの\R{点}{てん}と\R{青}{あお}の\R{点}{てん}を\R{結}{むす}ぶ\R{点線}{てん|せん}は、どれも\R{直線}{ちょく|せん} $y=x$（\R{灰色}{はい|いろ}の\R{点線}{てん|せん}）と\R{直角}{ちょっ|かく}に\R{交}{まじ}わり、そのまん\R{中}{なか}が $y=x$ の\R{上}{うえ}にあります。たとえば $(2,\ 4)$ と $(4,\ 2)$ のまん\R{中}{なか}は $(3,\ 3)$ です。

だから、\R{紙}{かみ}を $y=x$ で\R{折}{お}ると、2つのグラフはぴったり\R{重}{かさ}なります。
\end{minipage}

\par\vspace{4pt}
\begin{taisetsu}
$y=\log_ax$ のグラフは、$y=a^x$ のグラフを、\R{直線}{ちょく|せん} $y=x$ で\R{折}{お}り\R{返}{かえ}したもの。
\end{taisetsu}

グラフから、$y=\log_2x$ の\R{性質}{せい|しつ}が\R{読}{よ}みとれます。
\begin{enumerate}[label={\bkcircled[bkTheme]{\arabic*}}, leftmargin=2.4\zw, itemsep=1pt, topsep=2pt]
\item \R{点}{てん} $(1,\ 0)$ を\R{通}{とお}る（$\log_21=0$）。$y=2^x$ が $(0,\ 1)$ を\R{通}{とお}るのを\R{折}{お}り\R{返}{かえ}したもの。
\item $x$ は\R{正}{せい}の\R{数}{かず}のところにしかない（②の「0やマイナスの\R{対数}{たい|すう}はない」）。
\item \R{右上}{みぎ|あ}がりだが、とてもゆっくり\R{上}{あ}がる。$x=1024$ でも、$y$ はやっと10。
\end{enumerate}
\R{倍々}{ばい|ばい}で\R{急}{きゅう}にふえる\R{指数関数}{し|すう|かん|すう}を\R{折}{お}り\R{返}{かえ}したので、\R{対数}{たい|すう}は、とてもゆっくりふえる\R{関数}{かん|すう}になるのです。

\begin{motto}[3]
\par\addvspace{3pt}\textbf{\R{逆関数}{ぎゃく|かん|すう}}\par\noindent
$y=2^x$ は「\R{回数}{かい|すう} $x$ から\R{倍率}{ばい|りつ} $y$ を\R{出}{だ}す\R{箱}{はこ}」、$y=\log_2x$ は「\R{倍率}{ばい|りつ} $x$ から\R{回数}{かい|すう} $y$ を\R{出}{だ}す\R{箱}{はこ}」です（3-3の\R{入}{い}れると\R{出}{で}てくる\R{箱}{はこ}）。このように、\R{入}{い}れるものと\R{出}{で}てくるものを\R{入}{い}れかえた\R{関数}{かん|すう}を\textbf{\R{逆関数}{ぎゃく|かん|すう}}といいます。\R{逆関数}{ぎゃく|かん|すう}のグラフは、いつも\R{直線}{ちょく|せん} $y=x$ で\R{折}{お}り\R{返}{かえ}した\R{形}{かたち}です。5-14では、この\R{折}{お}り\R{返}{かえ}しを\R{使}{つか}って、\R{対数}{たい|すう}を\R{微分}{び|ぶん}します。
\end{motto}

\bksubsection{5}{\R{底}{てい}を\R{変}{か}える：\R{電卓}{でん|たく}で $\log_2$ を\R{求}{もと}める}

\leavevmode\label{bk:3-12:soko}\R{細菌}{さい|きん}が1000\R{倍}{ばい}になるのは、\R{何回}{なん|かい}\R{分}{わ}かれたときでしょう。\R{答}{こた}えは $\log_21000$ です。でも、\R{電卓}{でん|たく}の \texttt{log} のキーは、\R{底}{てい}が10の\R{常用対数}{じょう|よう|たい|すう}です。そこで、\R{常用対数}{じょう|よう|たい|すう}になおします。
\begin{enumerate}[label={\bkcircled[bkTheme]{\arabic*}}, leftmargin=2.4\zw, itemsep=1pt, topsep=2pt]
\item $\log_21000=p$ とすると、\R{指数}{し|すう}の\R{形}{かたち}で $2^p=1000$。
\item \R{両辺}{りょう|へん}の\R{常用対数}{じょう|よう|たい|すう}をとる（\R{等}{ひと}しい\R{数}{かず}は、\R{対数}{たい|すう}も\R{等}{ひと}しい）：$\log_{10}2^p=\log_{10}1000$。
\item \R{左辺}{さ|へん}は③の②で $p\log_{10}2$、\R{右辺}{う|へん}は3。だから $p\times0.301=3$。
\item \R{両辺}{りょう|へん}を0.301でわって（2-14）、$p=\dfrac{3}{0.301}=9.96\cdots$
\end{enumerate}
\R{約}{やく}9.97\R{回}{かい}です。\R{回数}{かい|すう}は\R{整数}{せい|すう}なので、10\R{回目}{かい|め}ではじめて1000\R{倍}{ばい}をこえます。たしかに、$2^9=512$、$2^{10}=1024$ です。

\R{同}{おな}じように\R{考}{かんが}えると、どんな\R{底}{てい}でも、\R{常用対数}{じょう|よう|たい|すう}のわり\R{算}{ざん}で\R{求}{もと}められます。
\bkthrmline{\fbox{$\log_ab=\dfrac{\log_{10}b}{\log_{10}a}$}}
\R{電卓}{でん|たく}には、\texttt{log} のほかに、もう1つ\R{対数}{たい|すう}のキーがあることがあります。それは、5-14で\R{学}{まな}ぶ「ふしぎな\R{数}{かず}」を\R{底}{てい}にした\R{対数}{たい|すう}です。

\bksubsection{6}{\R{倍々}{ばい|ばい}が\R{直線}{ちょく|せん}に\R{見}{み}える：\R{片対数}{かた|たい|すう}グラフ}

\leavevmode\label{bk:3-12:kata}\R{細菌}{さい|きん}の\R{数}{かず} $y=2^x$ を、$x=0$ から10までふつうのグラフにかくと、\R{図}{ず}\ref{bk:3-12:fig-futsu}のようになります。$y$ は1から1024まで\R{変}{か}わるので、はじめのほうの\R{点}{てん}は\R{横軸}{よこ|じく}にくっついて、ちがいが\R{読}{よ}めません。
\par\vspace{2pt}
\noindent\begin{minipage}[t]{0.49\linewidth}\centering
\begin{tikzpicture}
\begin{axis}[bkplot, width=60mm, height=58mm, xmin=0, xmax=11, ymin=0, ymax=1150, xtick={0,2,...,10}, ytick={0,200,...,1000}, xlabel={$x$}, ylabel={$y$}]
  \addplot[bkcurve, domain=0:10] {2^x};
  \addplot[only marks, mark=*, mark size=1.4pt, bkInk, samples at={0,1,...,10}] {2^x};
\end{axis}
\end{tikzpicture}
\bkfigcaption{ふつうのグラフ}\label{bk:3-12:fig-futsu}
\end{minipage}\hfill
\begin{minipage}[t]{0.49\linewidth}\centering
\begin{tikzpicture}
\begin{axis}[bkplot, width=60mm, height=58mm, xmin=0, xmax=11, ymin=0, ymax=3.25, xtick={0,2,...,10}, ytick={0,1,2,3}, yticklabels={1,10,100,1000},
  minor ytick={0.301,0.477,0.602,0.699,0.778,0.845,0.903,0.954,1.301,1.477,1.602,1.699,1.778,1.845,1.903,1.954,2.301,2.477,2.602,2.699,2.778,2.845,2.903,2.954,3.301}, yminorgrids=true, minor grid style={bkLine, line width=0.2pt}, minor tick length=0pt, xlabel={$x$}, ylabel={$y$}]
  \addplot[bkcurve, domain=0:10] {x*log10(2)};
  \addplot[only marks, mark=*, mark size=1.4pt, bkInk, samples at={0,1,...,10}] {x*log10(2)};
  \node[font=\sffamily\footnotesize, anchor=east] at (axis cs:-0.25,0) {1};
\end{axis}
\end{tikzpicture}
\bkfigcaption{\R{片対数}{かた|たい|すう}グラフ}\label{bk:3-12:fig-kata}
\end{minipage}

\par\vspace{4pt}
\R{図}{ず}\ref{bk:3-12:fig-kata}では、たて\R{軸}{じく}を\R{図}{ず}\ref{bk:3-12:fig-keta}の\R{対数目}{たい|すう|め}もりにしました。1、10、100、1000が\R{同}{おな}じ\R{間}{ま}かくで\R{並}{なら}び、そのあいだの\R{細}{ほそ}い\R{線}{せん}は2、3、…、9\R{倍}{ばい}の\R{所}{ところ}です。このようなグラフを\textbf{\R{片対数}{かた|たい|すう}グラフ}といいます。\bkidx{かたたいすうぐらふ}{\R{片対数}{かた|たい|すう}グラフ}\R{点}{てん}が、まっすぐに\R{並}{なら}びました。

\textbf{なぜ\R{直線}{ちょく|せん}になるのか}\quad たて\R{軸}{じく}の\R{位置}{い|ち}は、$y$ そのものではなく $\log_{10}y$ です。$y=2^x$ の\R{常用対数}{じょう|よう|たい|すう}をとると、③の②のきまりで
\[ \log_{10}y=\log_{10}2^x=x\times\log_{10}2=0.301x \]
$\log_{10}y$ は $x$ の1\R{次関数}{じ|かん|すう}（\R{比例}{ひ|れい}）なので、グラフは\R{傾}{かたむ}き0.301の\R{直線}{ちょく|せん}です（3-5）。\R{表}{ひょう}で\R{確}{たし}かめましょう。
\begin{bkcenter}
\begin{bktable}{colspec={Q[l,m,wd=22mm] *{7}{Q[c,m,wd=11.5mm]}}, rowsep=1.5pt, colsep=3pt, cells={font=\small}, row{3}={bg=bkthrmsel}}
$x$ & 0 & 1 & 2 & 3 & 4 & 5 & 10 \\
$y=2^x$ & 1 & 2 & 4 & 8 & 16 & 32 & 1024 \\
$\log_{10}y$ & 0 & 0.301 & 0.602 & 0.903 & 1.204 & 1.505 & 3.010 \\
\end{bktable}
\bktabcaption{$x$ が1ふえるごとに、$\log_{10}y$ は0.301ずつふえる（ひき\R{算}{ざん}が\R{一定}{いっ|てい}）。}
\end{bkcenter}
3-11で、\R{指数関数}{し|すう|かん|すう}は「わり\R{算}{ざん}が\R{一定}{いっ|てい}」、1\R{次関数}{じ|かん|すう}は「ひき\R{算}{ざん}が\R{一定}{いっ|てい}」でした。\R{対数}{たい|すう}をとると、かけ\R{算}{ざん}がたし\R{算}{ざん}になるので、「わり\R{算}{ざん}が\R{一定}{いっ|てい}」が「ひき\R{算}{ざん}が\R{一定}{いっ|てい}」に\R{変}{か}わるのです。

\textbf{\R{傾}{かたむ}きから\R{倍率}{ばい|りつ}がわかる}\quad \R{片対数}{かた|たい|すう}グラフで\R{直線}{ちょく|せん}の\R{傾}{かたむ}きが0.301なら、1つ\R{進}{すす}むごとの\R{倍率}{ばい|りつ}は $10^{0.301}=2$ \R{倍}{ばい}です（②の $a^{\log_ab}=b$）。はかったデータを\R{片対数}{かた|たい|すう}グラフにかいて\R{直線}{ちょく|せん}に\R{並}{なら}べば、その\R{量}{りょう}は\R{指数関数}{し|すう|かん|すう}で\R{変}{か}わっているとわかり、\R{倍率}{ばい|りつ}も\R{読}{よ}めます（\R{直線}{ちょく|せん}の\R{引}{ひ}き\R{方}{かた}は3-16）。

\bkthrmstep{1}{\R{両対数}{りょう|たい|すう}グラフ：\R{累乗}{るい|じょう}の\R{指数}{し|すう}が\R{傾}{かたむ}きになる}
\R{横軸}{よこ|じく}もたて\R{軸}{じく}も\R{対数目}{たい|すう|め}もりにしたものを\textbf{\R{両対数}{りょう|たい|すう}グラフ}といいます。\R{正方形}{せい|ほう|けい}の1\R{辺}{ぺん} $x$〔cm〕と\R{面積}{めん|せき} $y=x^2$\,\unit{cm^2} をかいてみましょう。
\par\vspace{2pt}
\noindent\begin{minipage}[c]{0.5\linewidth}\centering
\begin{bktable}{colspec={Q[l,m,wd=18mm] *{3}{Q[c,m,wd=13mm]}}, rowsep=1.5pt, colsep=3pt, cells={font=\small}}
$x$ & 1 & 10 & 100 \\
$y=x^2$ & 1 & 100 & 10000 \\
\SetCell{bg=bkthrmsel} $\log_{10}x$ & 0 & 1 & 2 \\
\SetCell{bg=bkthrmsel} $\log_{10}y$ & 0 & 2 & 4 \\
\end{bktable}
\par\vspace{4pt}
{\small $\log_{10}x$ が1ふえると、$\log_{10}y$ は2ふえる。}
\end{minipage}\hfill
\begin{minipage}[c]{0.46\linewidth}\centering
\begin{tikzpicture}
\begin{axis}[bkplot, width=54mm, height=54mm, xmin=0, xmax=2.3, ymin=0, ymax=4.4, xtick={1,2}, xticklabels={10,100}, ytick={1,2,3,4}, yticklabels={10,100,1000,10000}, xlabel={$x$}, ylabel={$y$}]
  \addplot[bkcurve, domain=0:2.1] {2*x};
  \addplot[only marks, mark=*, mark size=1.5pt, bkInk] coordinates {(0,0) (1,2) (2,4)};
  \node[font=\sffamily\footnotesize, anchor=north east] at (axis cs:-0.05,-0.05) {1};
  \node[font=\sffamily\footnotesize, text=bkBlue!80!black, anchor=west, fill=white, inner sep=1pt] at (axis cs:1.1,1.5) {$y=x^2$};
\end{axis}
\end{tikzpicture}
\end{minipage}
\par\vspace{4pt}
\R{点}{てん}は\R{傾}{かたむ}き2の\R{直線}{ちょく|せん}に\R{並}{なら}びます。$y=x^2$ の\R{両辺}{りょう|へん}の\R{常用対数}{じょう|よう|たい|すう}をとると、③の②のきまりで $\log_{10}y=2\log_{10}x$ になるからです。\R{傾}{かたむ}きの2は、$x^2$ の\R{指数}{し|すう}の2です。\R{振}{ふ}り\R{子}{こ}の\R{周期}{しゅう|き}は\R{糸}{いと}の\R{長}{なが}さの $\sqrt{\phantom{x}}$（$\dfrac{1}{2}$ \R{乗}{じょう}）に\R{比例}{ひ|れい}するので（2-8）、\R{両対数}{りょう|たい|すう}グラフでは\R{傾}{かたむ}き $\dfrac{1}{2}$ の\R{直線}{ちょく|せん}です。はかったデータの\R{傾}{かたむ}きを\R{読}{よ}めば、「\R{何乗}{なん|じょう}に\R{比例}{ひ|れい}するか」がわかります。

\bksubsection{7}{\R{比}{ひ}をけたで\R{表}{あらわ}す：デシベル}

\leavevmode\label{bk:3-12:db}\R{音}{おと}の\R{強}{つよ}さや\R{電気}{でん|き}の\R{信号}{しん|ごう}の\R{強}{つよ}さは、とても\R{広}{ひろ}いはんいで\R{変}{か}わります。\R{人}{ひと}の\R{耳}{みみ}は、ささやき\R{声}{ごえ}の\R{約}{やく}10\R{億}{おく}\R{倍}{ばい}の\R{強}{つよ}さの\R{音}{おと}まで\R{聞}{き}けます。そこで、2つの\R{強}{つよ}さの\R{比}{ひ}を、\R{常用対数}{じょう|よう|たい|すう}で「けた」として\R{表}{あらわ}します。
\begin{enumerate}[label={\bkcircled[bkTheme]{\arabic*}}, leftmargin=2.4\zw, itemsep=1pt, topsep=2pt]
\item 2つの\R{強}{つよ}さ（\R{電力}{でん|りょく}やエネルギー）の\R{比}{ひ}を\R{求}{もと}める。\R{例}{れい}：1000\R{倍}{ばい}。
\item その\R{常用対数}{じょう|よう|たい|すう}をとる。$\log_{10}1000=3$。これを3ベル（B）とよぶ。
\item ベルは\R{大}{おお}きすぎるので、10\R{倍}{ばい}した\R{数}{かず}を\R{使}{つか}う。$10\times3=30$。\R{単位}{たん|い}はデシベル（dB）。
\end{enumerate}
「デシ」は「10\R{分}{ぶん}の1」という\R{意味}{い|み}で、dL（デシリットル）のデシと\R{同}{おな}じです。1\,dB は $\dfrac{1}{10}$\,B なので、B の\R{数}{かず}を10\R{倍}{ばい}すると dB の\R{数}{かず}になります。\bkidx{でしべる}{デシベル}
\bkthrmline{\fbox{（デシベル）＝10×$\log_{10}$（\R{強}{つよ}さの\R{比}{ひ}）}}
\noindent\begin{minipage}[t]{0.42\linewidth}\vspace{0pt}
B（ベル）は、\R{電話}{でん|わ}を\R{発明}{はつ|めい}したベルの\R{名前}{な|まえ}からとった\R{単位}{たん|い}です。dB の d は、デシ（10\R{分}{ぶん}の1）の\R{頭文字}{かしら|も|じ}です。
\end{minipage}\hfill
\begin{minipage}[t]{67mm}\vspace{0pt}\bkkigocard{\sffamily dB}{デシベル}{\R{強}{つよ}さの\R{比}{ひ}を「10×\R{常用対数}{じょう|よう|たい|すう}」で\R{表}{あらわ}す\R{単位}{たん|い}。10\R{倍}{ばい}で10\,dB}{decibel}{3-12}\end{minipage}
\par\vspace{6pt}
\begin{bkcenter}
\begin{bktable}{colspec={Q[l,m,wd=28mm] *{6}{Q[c,m,wd=10.5mm]} Q[c,m,wd=18mm]}, rowsep=1.5pt, colsep=3pt, cells={font=\small}, row{2}={bg=bkthrmsel}}
\R{強}{つよ}さの\R{比}{ひ} & $\dfrac{1}{2}$ & 1 & 2 & 10 & 100 & 1000 & 1000000 \\
デシベル & $-3$ & 0 & 3 & 10 & 20 & 30 & 60 \\
\end{bktable}
\bktabcaption{2\R{倍}{ばい}は $10\times0.301=3.01$ で\R{約}{やく}3\,dB。$\dfrac{1}{2}$ は\R{約}{やく} $-3$\,dB。10\R{倍}{ばい}ごとに10\,dB ふえる。}
\end{bkcenter}

\bkhatenaref の\R{答}{こた}え：ふつうの\R{会話}{かい|わ}60\,dB とささやき\R{声}{ごえ}30\,dB の\R{差}{さ}は30\,dB です。$30\div10=3$ で、\R{常用対数}{じょう|よう|たい|すう}が3なので、\R{強}{つよ}さの\R{比}{ひ}は $10^3=1000$ \R{倍}{ばい}です。2\R{倍}{ばい}どころではありません。ガード\R{下}{した}の100\,dB は、ささやき\R{声}{ごえ}より70\,dB \R{大}{おお}きいので $10^7=1000$\R{万}{まん}\R{倍}{ばい}です。それでも\R{耳}{みみ}には1000\R{倍}{ばい}や1000\R{万}{まん}\R{倍}{ばい}もうるさくは\R{聞}{き}こえません。\R{人}{ひと}の\R{耳}{みみ}は、\R{強}{つよ}さの「けた」に\R{近}{ちか}い\R{感}{かん}じ\R{方}{かた}をするからです。デシベルは、\R{耳}{みみ}の\R{感}{かん}じ\R{方}{かた}にも\R{合}{あ}った\R{表}{あらわ}し\R{方}{かた}なのです。

\bkthrmstep{1}{つなぐと、たし\R{算}{ざん}}
\R{信号}{しん|ごう}を\R{大}{おお}きくする\R{装置}{そう|ち}（\R{増幅器}{ぞう|ふく|き}）を2つつなぎます。1つめが\R{電力}{でん|りょく}を100\R{倍}{ばい}（20\,dB）、2つめが1000\R{倍}{ばい}（30\,dB）にすると、\R{全体}{ぜん|たい}では $100\times1000=100000$ \R{倍}{ばい}です。dB では $20+30=50$\,dB で、$\log_{10}100000=5$、$10\times5=50$ と\R{合}{あ}います。かけ\R{算}{ざん}がたし\R{算}{ざん}になる（③の①）ので、\R{倍率}{ばい|りつ}をかけるかわりに、dB をたすだけですみます。
\begin{bkfigure}
\begin{tikzpicture}[x=1mm, y=1mm]
  \node[bk node, minimum width=12mm] (in) at (0,0) {\R{入}{い}り\R{口}{ぐち}};
  \node[bk node, minimum width=26mm, align=center] (a1) at (34,0) {×100\\20\,dB};
  \node[bk node, minimum width=26mm, align=center] (a2) at (72,0) {×1000\\30\,dB};
  \node[bk node, minimum width=12mm] (out) at (106,0) {\R{出}{で}\R{口}{ぐち}};
  \draw[bk link] (in) -- (a1); \draw[bk link] (a1) -- (a2); \draw[bk link] (a2) -- (out);
  \draw[decorate, decoration={brace, amplitude=3pt, mirror}, line width=0.6pt, bkOrange!80!black] (20,-9) -- (86,-9);
  \node[font=\sffamily\footnotesize, text=bkOrange!80!black, anchor=north] at (53,-12) {\R{全体}{ぜん|たい}：$100\times1000=100000$ \R{倍}{ばい}＝ $20+30=50$\,dB};
\end{tikzpicture}
\bkfigcaption{\R{倍率}{ばい|りつ}はかけ\R{算}{ざん}、dB はたし\R{算}{ざん}。}
\end{bkfigure}

\bkthrmstep{2}{\R{電圧}{でん|あつ}のときは20\R{倍}{ばい}する}
\R{抵抗}{てい|こう}に\R{流}{なが}れる\R{電力}{でん|りょく}は、\R{電圧}{でん|あつ}の2\R{乗}{じょう}に\R{比例}{ひ|れい}します（\R{電圧}{でん|あつ}が2\R{倍}{ばい}になると\R{電流}{でん|りゅう}も2\R{倍}{ばい}になり、\R{電力}{でん|りょく}は $2\times2=4$ \R{倍}{ばい}）。\R{電圧}{でん|あつ}の\R{比}{ひ}が $k$ なら、\R{電力}{でん|りょく}の\R{比}{ひ}は $k^2$ です。③の②のきまりで、
\[ 10\times\log_{10}k^2=10\times2\log_{10}k=20\log_{10}k \]
だから、\R{電圧}{でん|あつ}の\R{比}{ひ}から dB を\R{求}{もと}めるときは、\R{常用対数}{じょう|よう|たい|すう}を20\R{倍}{ばい}します。\R{電圧}{でん|あつ}が10\R{倍}{ばい}なら $20\times1=20$\,dB（\R{電力}{でん|りょく}は100\R{倍}{ばい}で、やはり20\,dB）です。

\begin{column}{\R{地震}{じ|しん}のマグニチュード}
\R{地震}{じ|しん}のエネルギーの\R{大}{おお}きさを\R{表}{あらわ}すマグニチュードも、\R{対数}{たい|すう}を\R{使}{つか}った\R{数}{かず}です。マグニチュードが1ふえると、エネルギーは $10^{1.5}=31.6\cdots$ で\R{約}{やく}32\R{倍}{ばい}、2ふえると $10^{1.5\times2}=10^3$ で1000\R{倍}{ばい}になります（3-11の\R{整数}{せい|すう}でない\R{指数}{し|すう}と、$(a^m)^n=a^{mn}$）。マグニチュード7の\R{地震}{じ|しん}は、マグニチュード5の\R{地震}{じ|しん}の\R{約}{やく}1000\R{倍}{ばい}のエネルギーです。
\end{column}

\begin{taisetsu}
$\log_ab$ は「$a$ を\R{何乗}{なん|じょう}すると $b$ になるか」。\par
\hspace*{1\zw}$a^p=b$ と $\log_ab=p$ は\R{同}{おな}じこと。\quad $a^{\log_ab}=b$\par\vspace{2pt}
$\log_aMN=\log_aM+\log_aN$、\ $\log_a\dfrac{M}{N}=\log_aM-\log_aN$\par\vspace{2pt}
$\log_aM^k=k\log_aM$\quad（$a>0$、$a\neq1$、$M>0$、$N>0$）\par\vspace{1pt}
{\mdseries\small \R{底}{てい}が10でも2でも、どんな\R{底}{てい} $a$ でも\R{成}{な}り\R{立}{た}ちます。どれも\R{指数}{し|すう}のきまり $a^p\times a^q=a^{p+q}$（3-11）から、\R{底}{てい}が\R{何}{なん}でも\R{同}{おな}じ\R{理由}{り|ゆう}で\R{出}{で}てくるからです（③の $b$、$c$ を $M$、$N$ と\R{書}{か}きました）。\par
\R{底}{てい}2で\R{確}{たし}かめ：$\log_28+\log_24=3+2=5=\log_232$（$8\times4=32$）\par}\vspace{2pt}
$\log_a1=0$、\ $\log_aa=1$\qquad $\log_ab=\dfrac{\log_{10}b}{\log_{10}a}$\par\vspace{2pt}
$y=\log_ax$ のグラフは、$y=a^x$ を\R{直線}{ちょく|せん} $y=x$ で\R{折}{お}り\R{返}{かえ}したもの。\par\vspace{2pt}
デシベル＝10×$\log_{10}$（\R{電力}{でん|りょく}の\R{比}{ひ}）＝20×$\log_{10}$（\R{電圧}{でん|あつ}の\R{比}{ひ}）
\end{taisetsu}

\begin{kotoba}
\small
\bkword{\R{対数}{たい|すう}・\R{底}{てい}}{$\log_ab$。$a$ を\R{何乗}{なん|じょう}すると $b$ になるかを\R{表}{あらわ}す\R{数}{かず}。$a$ が\R{底}{てい}。}
\bkword{\R{常用対数}{じょう|よう|たい|すう}}{\R{底}{てい}が10の\R{対数}{たい|すう}。けたの\R{数}{かず}がわかる。\R{電卓}{でん|たく}の \texttt{log}。}
\bkword{\R{片対数}{かた|たい|すう}グラフ}{たて\R{軸}{じく}だけを\R{対数目}{たい|すう|め}もりにしたグラフ。\R{指数関数}{し|すう|かん|すう}が\R{直線}{ちょく|せん}になる。}
\bkword{デシベル（dB）}{\R{強}{つよ}さの\R{比}{ひ}の\R{常用対数}{じょう|よう|たい|すう}を10\R{倍}{ばい}した\R{数}{かず}。}
\end{kotoba}

\begin{column}{pH（ピーエイチ）も\R{対数}{たい|すう}}
\R{水}{みず}よう\R{液}{えき}の\R{酸性}{さん|せい}の\R{強}{つよ}さを\R{表}{あらわ}す pH は、\R{水素}{すい|そ}イオンという\R{粒}{つぶ}の\R{濃}{こ}さの\R{常用対数}{じょう|よう|たい|すう}に、マイナスをつけた\R{数}{かず}です。\R{中性}{ちゅう|せい}の\R{水}{みず}は pH 7、レモン\R{汁}{じる}は\R{約}{やく} pH 2 です。\R{差}{さ}が5なので、レモン\R{汁}{じる}の\R{水素}{すい|そ}イオンの\R{濃}{こ}さは、\R{中性}{ちゅう|せい}の\R{水}{みず}の\R{約}{やく} $10^5=100000$ \R{倍}{ばい}です（マイナスがついているので、pH が\R{小}{ちい}さいほど\R{濃}{こ}い）。
\end{column}

%% ---------------- 例題 ----------------
\clearpage
\bksection{reidai}{\R{例題}{れい|だい}でたしかめよう}

\begin{reidai}[\R{指数}{し|すう}の\R{形}{かたち}になおして\R{対数}{たい|すう}を\R{求}{もと}める]
\R{次}{つぎ}の\R{値}{あたい}を\R{求}{もと}めましょう。\par
（1）\nobreak $\log_2 32$\quad（2）\nobreak $\log_3 81$\quad（3）\nobreak $\log_{10}0.001$\quad（4）\nobreak $\log_2\dfrac{1}{8}$\quad（5）\nobreak $\log_9 3$\quad（6）\nobreak $\log_5 1$\par
\kangaekata $\log_ab=\bkthrmbox$ を $a^{\bkthrmbox}=b$ になおし、$\bkthrmbox$ に\R{入}{はい}る\R{数}{かず}をさがします（②）。\par
\shiki（1）\nobreak $2^{\bkthrmbox}=32$。$2\times2\times2\times2\times2=32$ なので5\par
\hspace*{2.6\zw}（2）\nobreak $3^{\bkthrmbox}=81$。$3\times3\times3\times3=81$ なので4\par
\hspace*{2.6\zw}（3）\nobreak $10^{\bkthrmbox}=0.001=\dfrac{1}{1000}=\dfrac{1}{10^3}=10^{-3}$ なので $-3$\par
\hspace*{2.6\zw}（4）\nobreak $2^{\bkthrmbox}=\dfrac{1}{8}=\dfrac{1}{2^3}=2^{-3}$ なので $-3$\par
\hspace*{2.6\zw}（5）\nobreak $9^{\bkthrmbox}=3$。$9^{\frac{1}{2}}=\sqrt{9}=3$（3-11）なので $\dfrac{1}{2}$\par
\hspace*{2.6\zw}（6）\nobreak $5^{\bkthrmbox}=1$。$5^0=1$ なので0\par
\kotae（1）\nobreak 5\quad（2）\nobreak 4\quad（3）\nobreak $-3$\quad（4）\nobreak $-3$\quad（5）\nobreak $\dfrac{1}{2}$\quad（6）\nobreak 0
\end{reidai}

\begin{jibun}
\R{次}{つぎ}の\R{値}{あたい}を\R{求}{もと}めましょう。（1）\nobreak $\log_2 64$\quad（2）\nobreak $\log_3 27$\quad（3）\nobreak $\log_{10}10000$\quad（4）\nobreak $\log_2\dfrac{1}{4}$\quad（5）\nobreak $\log_4 2$\quad（6）\nobreak $\log_7 1$
\bkjans{（1）\nobreak 6（2）\nobreak 3（3）\nobreak 4（4）\nobreak $-2$（5）\nobreak $\dfrac{1}{2}$（6）\nobreak 0}
\end{jibun}

\begin{bkcenter}
\begin{bktable}{colspec={Q[l,m,wd=14mm] *{10}{Q[c,m,wd=9.4mm]}}, rowsep=1.2pt, colsep=1.5pt, cells={font=\small}, row{1}={bg=bkTheme!10}}
\R{指数}{し|すう} & 1 & 2 & 3 & 4 & 5 & 6 & 7 & 8 & 9 & 10 \\
$2^{\bkthrmbox}$ & 2 & 4 & 8 & 16 & 32 & 64 & 128 & 256 & 512 & 1024 \\
$3^{\bkthrmbox}$ & 3 & 9 & 27 & 81 & 243 & 729 & & & & \\
$5^{\bkthrmbox}$ & 5 & 25 & 125 & 625 & & & & & & \\
\end{bktable}
\bktabcaption{\R{累乗}{るい|じょう}の\R{早見表}{はや|み|ひょう}。\R{対数}{たい|すう}は、この\R{表}{ひょう}を\R{下}{した}から\R{上}{うえ}へ\R{読}{よ}めば\R{求}{もと}められる。}
\end{bkcenter}

\begin{reidai}[\R{対数}{たい|すう}のきまりで\R{常用対数}{じょう|よう|たい|すう}を\R{求}{もと}める]
$\log_{10}2=0.301$、$\log_{10}3=0.477$ を\R{使}{つか}って、\R{次}{つぎ}の\R{値}{あたい}を\R{求}{もと}めましょう。\par
（1）\nobreak $\log_{10}6$\quad（2）\nobreak $\log_{10}8$\quad（3）\nobreak $\log_{10}5$\quad（4）\nobreak $\log_{10}1.5$\quad（5）\nobreak $\log_{10}0.25$\par
\kangaekata 2と3と10のかけ\R{算}{ざん}・わり\R{算}{ざん}・\R{累乗}{るい|じょう}で\R{表}{あらわ}して、③のきまりを\R{使}{つか}います。$\log_{10}10=1$ です。\par
\shiki（1）\nobreak $6=2\times3$。かけ\R{算}{ざん}はたし\R{算}{ざん}になる\par
\hspace*{4.6\zw}$\log_{10}6=\log_{10}2+\log_{10}3=0.301+0.477=0.778$\par
\hspace*{2.6\zw}（2）\nobreak $8=2^3$。\R{累乗}{るい|じょう}はかけ\R{算}{ざん}になる\par
\hspace*{4.6\zw}$\log_{10}8=3\log_{10}2=3\times0.301=0.903$\par
\hspace*{2.6\zw}（3）\nobreak $5=\dfrac{10}{2}$。わり\R{算}{ざん}はひき\R{算}{ざん}になる\par
\hspace*{4.6\zw}$\log_{10}5=\log_{10}10-\log_{10}2=1-0.301=0.699$\par
\hspace*{2.6\zw}（4）\nobreak $1.5=\dfrac{3}{2}$。わり\R{算}{ざん}はひき\R{算}{ざん}\par
\hspace*{4.6\zw}$\log_{10}1.5=\log_{10}3-\log_{10}2=0.477-0.301=0.176$\par
\hspace*{2.6\zw}（5）\nobreak $0.25=\dfrac{1}{4}=\dfrac{1}{2^2}=2^{-2}$。\R{累乗}{るい|じょう}（$r=-2$）\par
\hspace*{4.6\zw}$\log_{10}0.25=-2\log_{10}2=-2\times0.301=-0.602$\par
\kotae（1）\nobreak 0.778\quad（2）\nobreak 0.903\quad（3）\nobreak 0.699\quad（4）\nobreak 0.176\quad（5）\nobreak $-0.602$
\end{reidai}

\begin{jibun}
$\log_{10}2=0.301$、$\log_{10}3=0.477$ を\R{使}{つか}って、\R{次}{つぎ}の\R{値}{あたい}を\R{求}{もと}めましょう。（1）\nobreak $\log_{10}12$\quad（2）\nobreak $\log_{10}9$\quad（3）\nobreak $\log_{10}20$\quad（4）\nobreak $\log_{10}0.5$
\bkjans{（1）\nobreak $12=2^2\times3$。$2\times0.301+0.477=1.079$（2）\nobreak $9=3^2$。$2\times0.477=0.954$（3）\nobreak $20=2\times10$。$0.301+1=1.301$（4）\nobreak $0.5=2^{-1}$。$-0.301$}
\end{jibun}
\noindent \R{答}{こた}えは\R{電卓}{でん|たく}の \texttt{log} のキーで\R{確}{たし}かめられます。たとえば $\log_{10}6=0.77815\cdots$、$\log_{10}5=0.69897\cdots$ で、\R{小数第}{しょう|すう|だい}3\R{位}{い}まで\R{合}{あ}っています。$\log_{10}2$ と $\log_{10}3$ さえわかれば、2と3と10のかけ\R{算}{ざん}・わり\R{算}{ざん}でできた\R{数}{かず}の\R{対数}{たい|すう}は、すべて\R{計算}{けい|さん}できるのです。

\begin{reidai}[けたの\R{数}{かず}を\R{求}{もと}める]
$\log_{10}2=0.301$、$\log_{10}5=0.699$ を\R{使}{つか}って、（1）$2^{30}$、（2）$5^{10}$ が\R{何}{なん}けたの\R{数}{かず}かを\R{求}{もと}めましょう。\par
\kangaekata \R{常用対数}{じょう|よう|たい|すう}を\R{求}{もと}め、その\R{整数}{せい|すう}の\R{部分}{ぶ|ぶん}に1をたします（①の②）。\par
\shiki（1）\nobreak $\log_{10}2^{30}=30\times0.301=9.03$\par
\hspace*{2.6\zw}9\R{以上}{い|じょう}10\R{未満}{み|まん}なので、$10^9\leqq2^{30}<10^{10}$。10けた\par
\hspace*{2.6\zw}（2）\nobreak $\log_{10}5^{10}=10\times0.699=6.99$\par
\hspace*{2.6\zw}6\R{以上}{い|じょう}7\R{未満}{み|まん}なので、$10^6\leqq5^{10}<10^7$。7けた\par
\kotae（1）\nobreak 10けた\quad（2）\nobreak 7けた
\end{reidai}
\noindent \R{確}{たし}かめ：$2^{30}=1073741824$（10けた）、$5^{10}=9765625$（7けた）です。3-11の\bkhatenaref の\R{細菌}{さい|きん}の\R{数}{かず}は、10けたの\R{数}{かず}だったのです。
\begin{bkfigure}
\begin{tikzpicture}[x=1mm, y=1mm]
  \draw[bkInk, line width=0.8pt] (0,0) -- (110,0);
  \foreach \k/\v in {0/8,1/9,2/10,3/11} {\draw[bkInk, line width=0.8pt] (\k*34+4,-1.6) -- (\k*34+4,1.6);
    \node[font=\small, anchor=north] at (\k*34+4,-2.4) {$\v$};
    \node[font=\small, anchor=south] at (\k*34+4,2.4) {$10^{\v}$};}
  \fill[bkRed] (38+0.03*34,0) circle (1.3);
  \draw[bkRed!80!black, line width=0.5pt] (38+0.03*34,1.4) -- (45,9);
  \node[font=\sffamily\footnotesize, text=bkRed!80!black, anchor=west] at (45.5,9.5) {$2^{30}$（$\log_{10}2^{30}=9.03$）};
  \node[font=\sffamily\footnotesize, anchor=east] at (-2,4) {\R{数}{かず}};
  \node[font=\sffamily\footnotesize, anchor=east] at (-2,-4) {\R{常用対数}{じょう|よう|たい|すう}};
  \fill[bkTheme!35] (38,-12) rectangle (72,-8.5);
  \node[font=\sffamily\footnotesize, anchor=north] at (55,-12.8) {10けたの\R{数}{かず}（$10^9$ \R{以上}{い|じょう} $10^{10}$ \R{未満}{み|まん}）};
\end{tikzpicture}
\bkfigcaption{\R{常用対数}{じょう|よう|たい|すう}が9と10のあいだなら、10けたの\R{数}{かず}。}
\end{bkfigure}

\begin{jibun}
$\log_{10}2=0.301$、$\log_{10}3=0.477$ を\R{使}{つか}って、（1）$3^{20}$、（2）$2^{50}$ が\R{何}{なん}けたの\R{数}{かず}かを\R{求}{もと}めましょう。
\bkjans{（1）\nobreak $20\times0.477=9.54$。$10^9\leqq3^{20}<10^{10}$ で10けた（2）\nobreak $50\times0.301=15.05$。$10^{15}\leqq2^{50}<10^{16}$ で16けた}
\end{jibun}

\begin{reidai}[\R{何回}{なん|かい}で\R{何倍}{なん|ばい}をこえるか]
（1）2\R{倍}{ばい}ずつふえる\R{細菌}{さい|きん}が、はじめて1000\R{倍}{ばい}をこえるのは\R{何回目}{なん|かい|め}ですか。\par
（2）1\R{年}{ねん}ごとに1.5\R{倍}{ばい}にふえる\R{量}{りょう}が、はじめて10\R{倍}{ばい}をこえるのは\R{何年後}{なん|ねん|ご}ですか。\par
（3）\R{半減期}{はん|げん|き}が6\R{時間}{じ|かん}の\R{物質}{ぶっ|しつ}（3-11）が、はじめの $\dfrac{1}{1000}$ \R{以下}{い|か}になるのは\R{何時間後}{なん|じ|かん|ご}ですか。\par
$\log_{10}2=0.301$、$\log_{10}1.5=0.176$（\bkreidairef{2}の（4））を\R{使}{つか}いましょう。\par
\kangaekata $a^n=k$ となる $n$ は $\log_ak=\dfrac{\log_{10}k}{\log_{10}a}$（⑤）。\R{回数}{かい|すう}は\R{整数}{せい|すう}なので、\R{小数}{しょう|すう}になったら\R{次}{つぎ}の\R{整数}{せい|すう}にします。\par
\shiki（1）\nobreak $\log_21000=\dfrac{\log_{10}1000}{\log_{10}2}=\dfrac{3}{0.301}=9.96\cdots$\par
\hspace*{2.6\zw}9\R{回}{かい}では $2^9=512$、10\R{回}{かい}で $2^{10}=1024$。10\R{回目}{かい|め}\par
\hspace*{2.6\zw}（2）\nobreak $\log_{1.5}10=\dfrac{\log_{10}10}{\log_{10}1.5}=\dfrac{1}{0.176}=5.68\cdots$\par
\hspace*{2.6\zw}5\R{年}{ねん}では $1.5^5=7.59\cdots$、6\R{年}{ねん}で $1.5^6=11.39\cdots$。6\R{年後}{ねん|ご}\par
\hspace*{2.6\zw}（3）\nobreak $n$ \R{回}{かい}\R{半分}{はん|ぶん}になると $\dfrac{1}{2^n}$。$\dfrac{1}{2^n}\leqq\dfrac{1}{1000}$ は $2^n\geqq1000$ と\R{同}{おな}じ\par
\hspace*{4.6\zw}（1）から $n=10$ \R{回}{かい}。$6\times10=60$\par
\kotae（1）\nobreak 10\R{回目}{かい|め}\quad（2）\nobreak 6\R{年後}{ねん|ご}\quad（3）\nobreak 60\R{時間後}{じ|かん|ご}
\end{reidai}

\begin{jibun}
（1）3\R{倍}{ばい}ずつふえる\R{量}{りょう}が、はじめて10000\R{倍}{ばい}をこえるのは\R{何回目}{なん|かい|め}ですか。（2）1\R{年}{ねん}ごとに1.2\R{倍}{ばい}にふえる\R{量}{りょう}が、はじめて2\R{倍}{ばい}をこえるのは\R{何年後}{なん|ねん|ご}ですか。（3）セシウム137（\R{半減期}{はん|げん|き}30\R{年}{ねん}）が、はじめの $\dfrac{1}{100}$ \R{以下}{い|か}になるのは\R{何年後}{なん|ねん|ご}ですか。$\log_{10}2=0.301$、$\log_{10}3=0.477$、$\log_{10}1.2=0.079$ を\R{使}{つか}いましょう。
\bkjans{（1）\nobreak $\dfrac{4}{0.477}=8.38\cdots$。$3^8=6561$、$3^9=19683$ で9\R{回目}{かい|め}（2）\nobreak $\dfrac{0.301}{0.079}=3.81\cdots$。$1.2^3=1.728$、$1.2^4=2.0736$ で4\R{年後}{ねん|ご}（3）\nobreak $2^n\geqq100$。$\dfrac{2}{0.301}=6.64\cdots$。$2^6=64$、$2^7=128$ で7\R{回}{かい}。$30\times7=210$ で210\R{年後}{ねん|ご}}
\end{jibun}

\begin{reidai}[デシベル]
（1）\R{電力}{でん|りょく}が1000\R{倍}{ばい}になるのは\R{何}{なん}dB ですか。（2）\R{電力}{でん|りょく}が $\dfrac{1}{2}$ になるのは\R{約}{やく}\R{何}{なん}dB ですか。（3）\R{電圧}{でん|あつ}が10\R{倍}{ばい}になるのは\R{何}{なん}dB ですか。（4）20\,dB と30\,dB の\R{増幅器}{ぞう|ふく|き}をつなぐと、\R{全体}{ぜん|たい}は\R{何}{なん}dB で、\R{電力}{でん|りょく}は\R{何倍}{なん|ばい}ですか。$\log_{10}2=0.301$ を\R{使}{つか}いましょう。\par
\kangaekata \R{電力}{でん|りょく}は「10×\R{常用対数}{じょう|よう|たい|すう}」、\R{電圧}{でん|あつ}は「20×\R{常用対数}{じょう|よう|たい|すう}」（⑦）。つなぐと dB はたし\R{算}{ざん}です。\par
\shiki（1）\nobreak $10\times\log_{10}1000=10\times3=30$\par
\hspace*{2.6\zw}（2）\nobreak $10\times\log_{10}\dfrac{1}{2}=10\times(-0.301)=-3.01$\par
\hspace*{2.6\zw}（3）\nobreak $20\times\log_{10}10=20\times1=20$\par
\hspace*{2.6\zw}（4）\nobreak $20+30=50$。$50\div10=5$ なので $10^5=100000$\par
\kotae（1）\nobreak 30\,dB\quad（2）\nobreak \R{約}{やく} $-3$\,dB\quad（3）\nobreak 20\,dB\quad（4）\nobreak 50\,dB、100000\R{倍}{ばい}
\end{reidai}

\begin{jibun}
（1）\R{電力}{でん|りょく}が100\R{倍}{ばい}になるのは\R{何}{なん}dB ですか。（2）\R{電力}{でん|りょく}が2\R{倍}{ばい}になるのは\R{約}{やく}\R{何}{なん}dB ですか。（3）\R{電圧}{でん|あつ}が100\R{倍}{ばい}になるのは\R{何}{なん}dB ですか。（4）10\,dB と20\,dB の\R{増幅器}{ぞう|ふく|き}をつなぐと、\R{全体}{ぜん|たい}は\R{何}{なん}dB で、\R{電力}{でん|りょく}は\R{何倍}{なん|ばい}ですか。
\bkjans{（1）\nobreak $10\times2=20$ で20\,dB（2）\nobreak $10\times0.301=3.01$ で\R{約}{やく}3\,dB（3）\nobreak $20\times2=40$ で40\,dB（4）\nobreak 30\,dB。$10^3=1000$ で1000\R{倍}{ばい}}
\end{jibun}

\begin{bkfigure}
\begin{tikzpicture}[x=1mm, y=1mm]
  \draw[bkInk, line width=0.8pt] (0,0) -- (108,0);
  \foreach \k in {0,...,6} {\draw[bkInk, line width=0.8pt] (\k*18,-1.6) -- (\k*18,1.6);
    \pgfmathtruncatemacro\dbv{\k*10}\node[font=\small, anchor=north] at (\k*18,-2.4) {\dbv};}
  \foreach \k/\v in {0/1,1/10,2/100,3/1000,4/{$10^4$},5/{$10^5$},6/{$10^6$}} {\node[font=\small, anchor=south] at (\k*18,2.4) {\v};}
  \foreach \m in {2,...,9} {\pgfmathsetmacro\xx{18*log10(\m)}\draw[bkBlue, line width=0.5pt] (\xx,0) -- (\xx,1.1);}
  \node[font=\sffamily\footnotesize, anchor=east] at (-3,4) {\R{電力}{でん|りょく}の\R{比}{ひ}};
  \node[font=\sffamily\footnotesize, anchor=east] at (-3,-4) {dB};
  \pgfmathsetmacro\xa{18*log10(2)}
  \draw[bkRed!80!black, line width=0.5pt] (\xa,-1) -- (\xa,-9);
  \node[font=\sffamily\scriptsize, text=bkRed!80!black, anchor=north] at (\xa,-9.3) {2\R{倍}{ばい}＝\R{約}{やく}3\,dB};
\end{tikzpicture}
\bkfigcaption{dB は、\R{電力}{でん|りょく}の\R{比}{ひ}を\R{対数目}{たい|すう|め}もりで\R{読}{よ}んだもの。10\R{倍}{ばい}ごとに10\,dB。}
\end{bkfigure}

\begin{reidai}[\R{片対数}{かた|たい|すう}グラフの\R{傾}{かたむ}きから\R{倍率}{ばい|りつ}を\R{読}{よ}む]
ある\R{量}{りょう} $y$ を、1\R{時間}{じ|かん}ごとにはかって\R{片対数}{かた|たい|すう}グラフにかくと、\R{点}{てん}が\R{直線}{ちょく|せん}に\R{並}{なら}び、$\log_{10}y$ は1\R{時間}{じ|かん}ごとに0.477ずつふえていました。（1）$y$ は1\R{時間}{じ|かん}ごとに\R{何倍}{なん|ばい}になっていますか。（2）3\R{時間}{じ|かん}では\R{何倍}{なん|ばい}ですか。$\log_{10}3=0.477$ を\R{使}{つか}いましょう。\par
\kangaekata $\log_{10}y$ が0.477ふえることは、$y$ が $10^{0.477}$ \R{倍}{ばい}になることです（③の①：かけ\R{算}{ざん}がたし\R{算}{ざん}）。$a^{\log_ab}=b$（②）を\R{使}{つか}います。\par
\shiki（1）\nobreak $10^{0.477}=10^{\log_{10}3}=3$\par
\hspace*{2.6\zw}（2）\nobreak 3\R{時間}{じ|かん}で $0.477\times3=1.431$ ふえる。$3\times3\times3=27$\par
\kotae（1）\nobreak 3\R{倍}{ばい}\quad（2）\nobreak 27\R{倍}{ばい}
\end{reidai}

\begin{jibun}
\R{片対数}{かた|たい|すう}グラフで、$\log_{10}y$ が1\R{日}{にち}ごとに0.699ずつふえる\R{直線}{ちょく|せん}になりました。（1）$y$ は1\R{日}{にち}ごとに\R{何倍}{なん|ばい}になっていますか。（2）2\R{日}{にち}では\R{何倍}{なん|ばい}ですか。$\log_{10}5=0.699$ を\R{使}{つか}いましょう。
\bkjans{（1）\nobreak $10^{0.699}=10^{\log_{10}5}=5$ で5\R{倍}{ばい}（2）\nobreak $5\times5=25$ で25\R{倍}{ばい}}
\end{jibun}

\noindent\begin{minipage}[c]{0.5\linewidth}\centering
\begin{tikzpicture}
\begin{axis}[bkplot, width=60mm, height=52mm, xmin=0, xmax=3.6, ymin=0, ymax=1.75, xtick={1,2,3}, ytick={0,0.477,0.954,1.431}, yticklabels={1,3,9,27},
  minor ytick={0.301,0.602,0.699,0.778,0.845,0.903,1.301,1.477,1.602,1.699}, yminorgrids=true, minor grid style={bkLine, line width=0.2pt}, minor tick length=0pt, xlabel={\R{時間}{じ|かん}}, ylabel={$y$}]
  \addplot[bkcurve, domain=0:3.2] {0.477*x};
  \addplot[only marks, mark=*, mark size=1.5pt, bkInk] coordinates {(0,0) (1,0.477) (2,0.954) (3,1.431)};
  \node[font=\sffamily\footnotesize, anchor=east] at (axis cs:-0.08,0) {1};
  \draw[bkOrange!80!black, line width=0.8pt] (axis cs:1,0.477) -- (axis cs:2,0.477) -- (axis cs:2,0.954);
  \node[font=\sffamily\scriptsize, text=bkOrange!80!black, anchor=north, fill=white, inner sep=1pt] at (axis cs:1.5,0.44) {1\R{時間}{じ|かん}};
  \node[font=\sffamily\scriptsize, text=bkOrange!80!black, anchor=west, fill=white, inner sep=1pt] at (axis cs:2.05,0.715) {＋0.477（×3）};
\end{axis}
\end{tikzpicture}
\end{minipage}\hfill
\begin{minipage}[c]{0.46\linewidth}
\bkreidairef{6}のようすを\R{片対数}{かた|たい|すう}グラフにかきました。たて\R{軸}{じく}の1、3、9、27は、\R{同}{おな}じ\R{間}{ま}かく（0.477ずつ）で\R{並}{なら}んでいます。1\R{時間}{じ|かん}\R{右}{みぎ}へ\R{進}{すす}むと、\R{直線}{ちょく|せん}は0.477だけ\R{上}{あ}がり、$y$ は3\R{倍}{ばい}になります。
\end{minipage}

\begin{reidai}[\R{指数関数}{し|すう|かん|すう}と\R{対数}{たい|すう}のグラフ]
（1）$y=2^x$ のグラフは\R{点}{てん} $(5,\ 32)$ を\R{通}{とお}ります。$y=\log_2x$ のグラフは、どの\R{点}{てん}を\R{通}{とお}りますか。（2）$y=\log_3x$ のグラフで、$x=9$ のときの $y$ を\R{求}{もと}めましょう。\par
\kangaekata $y=\log_ax$ のグラフは、$y=a^x$ のグラフを $y=x$ で\R{折}{お}り\R{返}{かえ}したもの。\R{点}{てん}の $x$ と $y$ が\R{入}{い}れかわります（④）。\par
\shiki（1）\nobreak $(5,\ 32)$ の $x$ と $y$ を\R{入}{い}れかえて $(32,\ 5)$。\R{確}{たし}かめ：$\log_232=5$\par
\hspace*{2.6\zw}（2）\nobreak $y=\log_39$。$3^2=9$ なので $y=2$\par
\kotae（1）\nobreak \R{点}{てん} $(32,\ 5)$\quad（2）\nobreak $y=2$
\end{reidai}

\begin{jibun}
（1）$y=3^x$ のグラフは\R{点}{てん} $(4,\ 81)$ を\R{通}{とお}ります。$y=\log_3x$ のグラフは、どの\R{点}{てん}を\R{通}{とお}りますか。（2）$y=\log_2x$ のグラフで、$x=\dfrac{1}{8}$ のときの $y$ を\R{求}{もと}めましょう。
\bkjans{（1）\nobreak \R{点}{てん} $(81,\ 4)$（2）\nobreak $2^{-3}=\dfrac{1}{8}$ なので $y=-3$}
\end{jibun}

%% ---------------- まちがえやすいところ ----------------
\Needspace{16\baselineskip}
\bksection{tsumazuki}{まちがえやすいところをチェック}\label{bk:3-12:tsuma}

\begin{tsumazuki}{たし\R{算}{ざん}の\R{対数}{たい|すう}を、\R{分}{わ}けてしまう}
\bkng{$\log(a+b)=\log a+\log b$}
\bkok{$\log(ab)=\log a+\log b$。$\log(a+b)$ は、それ\R{以上}{い|じょう}かんたんにならない}
たとえば $\log_{10}(1+1)=\log_{10}2=0.301$ ですが、$\log_{10}1+\log_{10}1=0+0=0$ です。\R{和}{わ}になるのは、かけ\R{算}{ざん}のときだけです。
\end{tsumazuki}

\begin{tsumazuki}{\R{対数}{たい|すう}を、わり\R{算}{ざん}とまちがえる}
\bkng{$\log_28=8\div2=4$}
\bkok{$2^3=8$ なので $\log_28=3$}
$\log_ab$ は「$a$ を\R{何乗}{なん|じょう}すると $b$ か」です。\R{指数}{し|すう}の\R{形}{かたち} $2^{\bkthrmbox}=8$ になおして\R{確}{たし}かめましょう。
\end{tsumazuki}

\begin{tsumazuki}{\R{累乗}{るい|じょう}の\R{対数}{たい|すう}で、\R{対数}{たい|すう}を\R{累乗}{るい|じょう}する}
\bkng{$\log_{10}2^3=(\log_{10}2)^3=0.301^3$}
\bkok{$\log_{10}2^3=3\log_{10}2=3\times0.301=0.903$}
$2^3=8$ で、$\log_{10}8=0.903$ です。$0.301^3$ は\R{約}{やく}0.027で、まったくちがいます。
\end{tsumazuki}

\begin{tsumazuki}{デシベルを、そのままたす}
\bkng{60\,dB のスピーカー2\R{台}{だい}を\R{同}{おな}じ\R{強}{つよ}さで\R{鳴}{な}らすと、$60+60=120$\,dB}
\bkok{\R{強}{つよ}さが2\R{倍}{ばい}になるので、$60+3=63$\,dB}
\R{強}{つよ}さの\R{比}{ひ}が2\R{倍}{ばい}なら、$10\times0.301$ で3\,dB ふえるだけです。dB をたしてよいのは、\R{増幅器}{ぞう|ふく|き}のように\R{倍率}{ばい|りつ}をかけるときです。
\end{tsumazuki}

\begin{tsumazuki}{0やマイナスの\R{数}{かず}の\R{対数}{たい|すう}}
\bkng{$\log_20=0$}
\bkok{$\log_20$ はない（2を\R{何乗}{なん|じょう}しても0にならない）}
\R{対数}{たい|すう}がゼロになるのは、$\log_a1=0$ です（$a^0=1$）。$\log_a$ の\R{中}{なか}は、いつも\R{正}{せい}の\R{数}{かず}です。
\end{tsumazuki}

\begin{tsumazuki}{\R{逆数}{ぎゃく|すう}の\R{対数}{たい|すう}を、\R{対数}{たい|すう}の\R{逆数}{ぎゃく|すう}にする}
\bkng{$\log_{10}\dfrac{1}{2}=\dfrac{1}{\log_{10}2}=\dfrac{1}{0.301}=3.32\cdots$}
\bkok{$\log_{10}\dfrac{1}{2}=\log_{10}2^{-1}=-\log_{10}2=-0.301$}
$\dfrac{1}{2}$ は1より\R{小}{ちい}さいので、10の\R{指数}{し|すう}はマイナスです（$10^0=1$）。3.32だと $10^{3.32}$ で2000より\R{大}{おお}きな\R{数}{かず}になってしまいます。
\end{tsumazuki}

\begin{tsumazuki}{わり\R{算}{ざん}の\R{対数}{たい|すう}を、\R{対数}{たい|すう}のわり\R{算}{ざん}にする}
\bkng{$\log_{10}\dfrac{6}{2}=\dfrac{\log_{10}6}{\log_{10}2}=\dfrac{0.778}{0.301}=2.58\cdots$}
\bkok{$\log_{10}\dfrac{6}{2}=\log_{10}6-\log_{10}2=0.778-0.301=0.477$}
$\dfrac{6}{2}=3$ なので、\R{答}{こた}えは $\log_{10}3=0.477$ のはずです。\R{対数}{たい|すう}では、わり\R{算}{ざん}はひき\R{算}{ざん}になります。$\dfrac{\log_{10}b}{\log_{10}a}$ の\R{形}{かたち}は、\R{底}{てい}を\R{変}{か}えるとき（⑤）の\R{式}{しき}です。
\end{tsumazuki}

\begin{tsumazuki}{\R{対数目}{たい|すう|め}もりを、ふつうの\R{目}{め}もりのように\R{読}{よ}む}
\bkng{\R{片対数}{かた|たい|すう}グラフで、1と10の\R{目}{め}もりのまん\R{中}{なか}は5.5}
\bkok{まん\R{中}{なか}は $10^{0.5}=\sqrt{10}=3.16\cdots$}
\R{対数目}{たい|すう|め}もりでは、\R{位置}{い|ち}が $\log_{10}y$ を\R{表}{あらわ}します。まん\R{中}{なか}は $\log_{10}y=0.5$ の\R{所}{ところ}です。5の\R{目}{め}もり（0.699）は、まん\R{中}{なか}より\R{右}{みぎ}（\R{上}{うえ}）にあります。
\end{tsumazuki}

%% ---------------- 物理・工学 ----------------
\begin{butsuri}
\par\addvspace{3pt}\textbf{ダイオード：\R{片対数}{かた|たい|すう}グラフで\R{直線}{ちょく|せん}になる\R{電流}{でん|りゅう}}\par\noindent
\begin{minipage}[c]{0.5\linewidth}
ダイオードは、\R{電流}{でん|りゅう}を1\R{方向}{ほう|こう}にしか\R{流}{なが}さない\R{半導体}{はん|どう|たい}の\R{部品}{ぶ|ひん}です（LED もなかま）。\R{電圧}{でん|あつ}を\R{少}{すこ}しずつ\R{上}{あ}げると、\R{電流}{でん|りゅう}は\R{電圧}{でん|あつ}が\R{約}{やく}0.06\,V（60\,mV）\R{上}{あ}がるごとに\R{約}{やく}10\R{倍}{ばい}になります（\R{室温}{しつ|おん}で、\R{理想}{り|そう}に\R{近}{ちか}いダイオードのとき）。\R{指数関数}{し|すう|かん|すう}なので、ふつうのグラフでは\R{急}{きゅう}に\R{立}{た}ち\R{上}{あ}がって\R{見}{み}えますが、\R{片対数}{かた|たい|すう}グラフでは\R{直線}{ちょく|せん}です。\R{半導体}{はん|どう|たい}の\R{設計}{せっ|けい}では、この\R{直線}{ちょく|せん}の\R{傾}{かたむ}きを\R{読}{よ}んで、\R{部品}{ぶ|ひん}のよしあしを\R{調}{しら}べます。
\end{minipage}\hfill
\begin{minipage}[c]{0.46\linewidth}\centering
\begin{tikzpicture}
\begin{axis}[bkplot, width=56mm, height=54mm, xmin=0.4, xmax=0.68, ymin=0, ymax=4.3, xtick={0.42,0.48,0.54,0.60,0.66}, xticklabel style={/pgf/number format/fixed, /pgf/number format/precision=2}, ytick={0,1,2,3,4}, yticklabels={0.001,0.01,0.1,1,10},
  xlabel={V}, ylabel={mA}, axis line style={-}, clip=true]
  \addplot[bkcurve, domain=0.41:0.67] {(x-0.42)/0.06};
  \addplot[only marks, mark=*, mark size=1.4pt, bkInk] coordinates {(0.42,0) (0.48,1) (0.54,2) (0.60,3) (0.66,4)};
\end{axis}
\end{tikzpicture}
\par{\sffamily\scriptsize 0.06\,V ごとに\R{電流}{でん|りゅう}が10\R{倍}{ばい}（\R{片対数}{かた|たい|すう}）}
\end{minipage}

\par\addvspace{6pt}\textbf{ボード\R{線図}{せん|ず}：\R{機械}{き|かい}や\R{回路}{かい|ろ}の「\R{入}{い}れると\R{出}{で}てくる」をdBで\R{見}{み}る}\par\noindent
スピーカーやモーターの\R{制御}{せい|ぎょ}では、ゆっくりから\R{速}{はや}くまで、いろいろな\R{速}{はや}さでゆれる\R{信号}{しん|ごう}を\R{入}{い}れて、\R{出}{で}てくる\R{信号}{しん|ごう}が\R{何倍}{なん|ばい}になるかを\R{調}{しら}べます。\R{横軸}{よこ|じく}にゆれの\R{速}{はや}さ（1\R{秒}{びょう}あたりのゆれの\R{回数}{かい|すう}）を\R{対数目}{たい|すう|め}もりで、たて\R{軸}{じく}に\R{倍率}{ばい|りつ}を dB でとったグラフを\textbf{ボード\R{線図}{せん|ず}}といいます。どちらの\R{軸}{じく}も「けた」で\R{見}{み}ているので、100\R{倍}{ばい}も1000\R{倍}{ばい}もちがう\R{範囲}{はん|い}が1\R{枚}{まい}のグラフにおさまり、\R{装置}{そう|ち}をつないだときの\R{倍率}{ばい|りつ}も、dB のたし\R{算}{ざん}で\R{求}{もと}められます（9-8）。
\par\vspace{3pt}
\noindent\begin{minipage}[c]{0.5\linewidth}\centering
\begin{tikzpicture}
\begin{axis}[bkplot, axis lines=left, width=60mm, height=44mm, xmin=1, xmax=4.25, ymin=-45, ymax=6, xtick={1,2,3,4}, xticklabels={10,100,1000,10000}, ytick={-40,-20,0}, xlabel={Hz}, ylabel={dB},
  every axis x label/.style={at={(ticklabel* cs:1)}, anchor=north west, font=\sffamily\small}]
  \addplot[bkcurve, domain=1:4, samples=80] {-10*log10(1+10^(2*(x-2)))};
\end{axis}
\end{tikzpicture}
\end{minipage}\hfill
\begin{minipage}[c]{0.46\linewidth}
\R{抵抗}{てい|こう}とコンデンサーで\R{作}{つく}る、\R{低}{ひく}い\R{音}{おと}だけを\R{通}{とお}す\R{回路}{かい|ろ}のボード\R{線図}{せん|ず}です。100\,Hz（1\R{秒}{びょう}に100\R{回}{かい}のゆれ）より\R{高}{たか}いと、ゆれの\R{回数}{かい|すう}が10\R{倍}{ばい}になるごとに20\,dB（\R{電圧}{でん|あつ}は $\dfrac{1}{10}$）ずつ\R{下}{さ}がり、\R{直線}{ちょく|せん}になります。
\end{minipage}
\end{butsuri}



%% ---------------- 章末問題 ----------------
\bkendproblems
\bkproblemgroup{\R{基本}{き|ほん}}{1}

\bkproblem \R{次}{つぎ}の\R{値}{あたい}を\R{求}{もと}めましょう。（1）\nobreak $\log_2 16$\quad（2）\nobreak $\log_5 125$\quad（3）\nobreak $\log_{10}0.1$\quad（4）\nobreak $\log_3\dfrac{1}{9}$\quad（5）\nobreak $\log_{16}4$\quad（6）\nobreak $\log_{10}1$
\bkpans{（1）\nobreak 4（2）\nobreak 3（3）\nobreak $-1$（4）\nobreak $-2$（5）\nobreak $\dfrac{1}{2}$（6）\nobreak 0}

\bkproblem $\log_{10}2=0.301$、$\log_{10}3=0.477$ を\R{使}{つか}って、\R{次}{つぎ}の\R{値}{あたい}を\R{求}{もと}めましょう。（1）\nobreak $\log_{10}18$\quad（2）\nobreak $\log_{10}16$\quad（3）\nobreak $\log_{10}50$\quad（4）\nobreak $\log_{10}\dfrac{2}{3}$
\bkpans{（1）\nobreak $18=2\times3^2$。$0.301+0.954=1.255$（2）\nobreak $16=2^4$。$4\times0.301=1.204$（3）\nobreak $50=\dfrac{100}{2}$。$2-0.301=1.699$（4）\nobreak $0.301-0.477=-0.176$}

\bkproblem $\log_{10}2=0.301$、$\log_{10}3=0.477$ を\R{使}{つか}って、（1）$2^{40}$、（2）$3^{30}$ が\R{何}{なん}けたの\R{数}{かず}かを\R{求}{もと}めましょう。
\bkpans{（1）\nobreak $40\times0.301=12.04$ で13けた（2）\nobreak $30\times0.477=14.31$ で15けた}

\bkproblem 2\R{倍}{ばい}ずつふえる\R{量}{りょう}が、はじめて10\R{万}{まん}\R{倍}{ばい}をこえるのは\R{何回目}{なん|かい|め}ですか。$\log_{10}2=0.301$ を\R{使}{つか}いましょう。
\bkpans{$\log_2100000=\dfrac{5}{0.301}=16.6\cdots$。$2^{16}=65536$、$2^{17}=131072$ で17\R{回目}{かい|め}}

\bkproblem （1）\R{電力}{でん|りょく}が10000\R{倍}{ばい}になるのは\R{何}{なん}dB ですか。（2）\R{電圧}{でん|あつ}が1000\R{倍}{ばい}になるのは\R{何}{なん}dB ですか。（3）\R{電力}{でん|りょく}が $\dfrac{1}{2}$ になる\R{部品}{ぶ|ひん}を2つつなぐと、\R{全体}{ぜん|たい}は\R{約}{やく}\R{何}{なん}dB ですか。$\log_{10}2=0.301$ を\R{使}{つか}いましょう。
\bkpans{（1）\nobreak $10\times4=40$ で40\,dB（2）\nobreak $20\times3=60$ で60\,dB（3）\nobreak $-3+(-3)=-6$ で\R{約}{やく} $-6$\,dB（\R{電力}{でん|りょく}は $\dfrac{1}{4}$）}

\bkproblem （1）$y=2^x$ のグラフは\R{点}{てん} $(6,\ 64)$ を\R{通}{とお}ります。$y=\log_2x$ のグラフは、どの\R{点}{てん}を\R{通}{とお}りますか。（2）\R{片対数}{かた|たい|すう}グラフで、$\log_{10}y$ が1\R{分}{ぷん}ごとに0.301ずつふえる\R{直線}{ちょく|せん}になりました。$y$ は1\R{分}{ぷん}ごとに\R{何倍}{なん|ばい}ですか。
\bkpans{（1）\nobreak \R{点}{てん} $(64,\ 6)$（2）\nobreak $10^{0.301}=2$ で2\R{倍}{ばい}}

\bkproblemgroup{チャレンジ}{2}

\bkproblem $2^{100}$ は\R{何}{なん}けたの\R{数}{かず}ですか。$\log_{10}2=0.301$ を\R{使}{つか}いましょう。
\bkpans{$100\times0.301=30.1$。$10^{30}\leqq2^{100}<10^{31}$ で31けた}

\bkproblem \R{工場}{こう|じょう}の\R{機械}{き|かい}の\R{音}{おと}が90\,dB、\R{静}{しず}かな\R{図書館}{と|しょ|かん}が40\,dB でした。\R{音}{おと}の\R{強}{つよ}さは\R{何倍}{なん|ばい}ちがいますか。
\bkpans{$90-40=50$\,dB。$50\div10=5$ なので $10^5=100000$ \R{倍}{ばい}}

\bkproblem \R{電圧}{でん|あつ}を10\R{倍}{ばい}にする\R{増幅器}{ぞう|ふく|き}と、\R{電圧}{でん|あつ}を100\R{倍}{ばい}にする\R{増幅器}{ぞう|ふく|き}をつなぎました。\R{全体}{ぜん|たい}は\R{何}{なん}dB ですか。また、\R{入力}{にゅう|りょく}と\R{出力}{しゅつ|りょく}が\R{同}{おな}じ\R{抵抗}{てい|こう}のとき、\R{電力}{でん|りょく}は\R{何倍}{なん|ばい}になりますか。
\bkpans{$20+40=60$\,dB。$60\div10=6$ なので\R{電力}{でん|りょく}は $10^6=1000000$ \R{倍}{ばい}（\R{電圧}{でん|あつ}は1000\R{倍}{ばい}）}

\bkproblem $\log_{10}2=0.301$、$\log_{10}3=0.477$ を\R{使}{つか}って、$\log_23$ を\R{小数第}{しょう|すう|だい}2\R{位}{い}まで\R{求}{もと}めましょう。また、その\R{意味}{い|み}を\R{言}{い}いましょう。
\bkpans{$\log_23=\dfrac{\log_{10}3}{\log_{10}2}=\dfrac{0.477}{0.301}=1.584\cdots$ で\R{約}{やく}1.58。2を\R{約}{やく}1.58\R{乗}{じょう}すると3になる}

\begin{furikaeri}
\begin{checklist}
\item 10を\R{何乗}{なん|じょう}するかでけたを\R{数}{かぞ}え、$\log_ab$ の\R{意味}{い|み}を\R{説明}{せつ|めい}できる。\bkfurref{bk:3-12:log}
\item \R{対数}{たい|すう}のきまりを、\R{指数}{し|すう}のきまりから\R{作}{つく}って\R{使}{つか}える。\bkfurref{bk:3-12:kimari}
\item \R{片対数}{かた|たい|すう}グラフやデシベルを\R{使}{つか}える。\bkfurref{bk:3-12:kata}
\end{checklist}
\end{furikaeri}

\bkfukidashi{delta}{かけ\R{算}{ざん}をたし\R{算}{ざん}に、\R{倍々}{ばい|ばい}を\R{直線}{ちょく|せん}に。\R{対数}{たい|すう}は、\R{大}{おお}きさが\R{何}{なん}けたもちがうものを、1つのものさしで\R{見}{み}る\R{道具}{どう|ぐ}だね。\R{次}{つぎ}の\R{章}{しょう}からは、データのばらつきを\R{調}{しら}べるよ。}
